\documentclass[11pt]{article}
\usepackage[margin=1in]{geometry}
\usepackage{amsmath,amssymb,amsthm,mathtools}
\usepackage{array,booktabs,longtable}
\usepackage[hidelinks]{hyperref}

\newtheorem{theorem}{Theorem}[section]
\newtheorem{proposition}[theorem]{Proposition}
\newtheorem{lemma}[theorem]{Lemma}
\theoremstyle{remark}\newtheorem{remark}[theorem]{Remark}
\newcommand{\SpinSeven}{\operatorname{Spin}(7)}
\newcommand{\Hol}{\operatorname{Hol}}
\newcommand{\Fix}{\operatorname{Fix}}
\newcommand{\ii}{\mathrm i}
\newcommand{\EH}{\mathrm{EH}}
\newcommand{\T}{\mathbb T}
\newcommand{\R}{\mathbb R}
\newcommand{\C}{\mathbb C}
\newcommand{\Z}{\mathbb Z}
\newcommand{\RP}{\mathbb {RP}}
\newcommand{\CP}{\mathbb {CP}}
\newcommand{\PD}{\operatorname{PD}}
\newcommand{\id}{\operatorname{id}}
\newcommand{\diag}{\operatorname{diag}}

\title{A nonformal compact manifold with full holonomy $\SpinSeven$\\
obtained from the Mart\'in--Merch\'an $G_2$ construction}
\author{}\date{}
\begin{document}
\maketitle

\begin{abstract}
We start from the flat orbifold underlying the nonformal $G_2$ manifold of
Mart\'in--Merch\'an and introduce an affine anti-$G_2$ involution whose
translation term is a real quarter-period in the third elliptic coordinate.
After adjoining a circle, we couple this involution to reflection in the
circle; the resulting involution preserves the product Cayley form.  We
compute the complete isotropy stratification of the resulting
eight-orbifold.  In addition to codimension-four $A_1$ strata, it has
codimension-six strata modelled on
$\R^2\times(\C^3/\Z_2^2)$ and isolated strata modelled on
$(\R^4/\{\pm1\})^2$; no isolated $\R^8/\{\pm1\}$ local model occurs.
Explicit Eguchi--Hanson, toric crepant, and product models satisfy all of
Joyce's $R$-data compatibility conditions and yield a compact torsion-free
$\SpinSeven$ resolution.  On this resolution we construct a defined triple
Massey product of degrees $(3,2,2)$ whose full Massey coset does not contain
zero, and we prove that the resolution is simply connected.  The
Berger--Simons--Wang classification then reduces every proper holonomy
alternative to a K\"ahler case, which nonformality excludes.  Therefore the
resulting metrics have full holonomy $\SpinSeven$.
\end{abstract}

All cohomology has real coefficients unless otherwise stated.  Products of
transformations are composed from right to left.

\section{Statement of the construction}
We set
\[
E=\C/(\Z+\ii\Z),\qquad \T^7=(\R/4\Z)_t\times E^3,
\]
and write $z_j=x_j+\ii y_j$.  We define
\begin{align}
F(t,z_1,z_2,z_3)&=(t+1,\ii z_1,\ii z_2,-z_3),\label{eq:F}\\
\iota _1(t,z_1,z_2,z_3)&=(t,-z_1,-z_2,z_3+\tfrac12),\label{eq:i1}\\
\iota _2(t,z_1,z_2,z_3)&=(t,-z_1,-z_2,z_3+\tfrac{\ii}{2}),\label{eq:i2}\\
\kappa(t,z_1,z_2,z_3)&=(1-t,\bar z_1,\bar z_2,\bar z_3).\label{eq:kappa}
\end{align}
They preserve
\begin{equation}\label{eq:g2form}
\varphi=dt\wedge\omega+\operatorname{Re}\Theta,
\qquad \omega=\sum_{j=1}^3dx_j\wedge dy_j,
\qquad \Theta=dz_1\wedge dz_2\wedge dz_3.
\end{equation}
The group $G=\langle F,\iota _1,\iota _2,\kappa\rangle
\cong D_4\times\Z_2^2$ has order $32$, and $X=\T^7/G$ is the flat $G_2$
orbifold of \cite{MM}.

We define the additional affine involution
\begin{equation}\label{eq:tau}
\tau(t,z_1,z_2,z_3)=(1-t,z_1,z_2,-z_3+\tfrac14).
\end{equation}
The translation term $1/4$ in the third complex coordinate is one quarter
of the real period of $E$.  We set $S^1_s=\R/\Z$ and define
\begin{equation}\label{eq:cayley-product}
\Phi_0=ds\wedge\varphi+*_{\varphi}\varphi
\quad\text{on}\quad \T^8=S^1_s\times\T^7.
\end{equation}
We extend $G$ trivially in the $s$-direction and define
\begin{equation}\label{eq:T-def}
T(s,p)=(-s,\tau p),\qquad K=\langle G,T\rangle,
\qquad \mathcal O=\T^8/K.
\end{equation}

\begin{theorem}\label{thm:main}
There exist a compact smooth eight-manifold $\widehat Z$, a proper map
$\pi:\widehat Z\to\mathcal O$ that is a diffeomorphism over the regular
locus, and a number $\varepsilon>0$ such that, for every
$t\in(0,\varepsilon)$, $\widehat Z$ carries a torsion-free $\SpinSeven$
structure $\widehat\Phi_t$ whose associated metric has full holonomy
$\SpinSeven$.  Moreover, $\pi_1(\widehat Z)=0$ and $\widehat Z$ is
nonformal.
\end{theorem}

\begin{remark}\label{rem:not-product}
The product $S^1\times M$ of the Mart\'in--Merch\'an manifold with a circle
has a torsion-free $\SpinSeven$ form, but the product metric has holonomy
contained in $G_2$ because its circle vector field is parallel.  The
construction instead couples the reflection $s\mapsto-s$ with the
involution $\tau$, which reverses the $G_2$ form.  Section~5 proves that the
resulting resolution is nonformal.  Section~6 proves that it is simply
connected and then applies the de~Rham and Berger--Simons--Wang
classifications to show that the holonomy of each sufficiently small
torsion-free structure is exactly $\SpinSeven$.
\end{remark}

\section{The affine group and its fixed sets}
\subsection{Normalization and the Cayley form}

\begin{lemma}\label{lem:normalization}
The map $\tau$ normalizes $G$, with
\begin{equation}\label{eq:tau-relations}
\tau^2=1,\quad \tau F\tau=F^{-1}\iota _1,\quad
\tau\iota_j\tau=\iota_j\ (j=1,2),\quad \tau\kappa\tau=\kappa.
\end{equation}
Moreover,
\begin{equation}\label{eq:tau-forms}
\tau^*\varphi=-\varphi,\qquad
\tau^*(*_{\varphi}\varphi)=*_{\varphi}\varphi.
\end{equation}
Consequently, $|K|=64$ and every element of $K$ preserves $\Phi_0$.
\end{lemma}
\begin{proof}
Substituting the formulas \eqref{eq:F}--\eqref{eq:kappa} and
\eqref{eq:tau} into each composite gives all four relations in
\eqref{eq:tau-relations}.  The derivative of $\tau$ sends $dt$ and $dz_3$
to their negatives and fixes $dz_1,dz_2$; hence it preserves every summand
of $\omega$ and sends $\Theta$ to $-\Theta$.  This proves the first equation in
\eqref{eq:tau-forms}; the second follows from
$*\varphi=\frac12\omega^2-dt\wedge\operatorname{Im}\Theta$.  Since
$T^*ds=-ds$, one has $T^*\Phi_0=\Phi_0$.  Finally $T\notin G$ and normalizes
$G$, so $|K|=64$.
\end{proof}

\subsection{The coset \texorpdfstring{$GT$}{GT}}
We set $\bar t=t/4$ and use the coordinates
\begin{equation}\label{eq:seven-coordinates}
q=(\bar t,x_1,y_1,x_2,y_2,x_3,y_3)\in(\R/\Z)^7.
\end{equation}
For $q\mapsto Aq+v$, the fixed equation is
\begin{equation}\label{eq:fixed-equation}(A-I)q=-v\pmod{\Z^7}.\end{equation}

\begin{lemma}\label{lem:torus-congruence}
Let $M\in M_{m\times n}(\Z)$ and $y\in\R^m$.  The congruence
\[
 Mx=y\pmod{\Z^m},\qquad x\in\R^n/\Z^n,
\]
is solvable if and only if
\begin{equation}\label{eq:left-kernel-test}
 \lambda^Ty\in\Z
 \quad\text{for every }\lambda\in\ker(M^T)\cap\Z^m.
\end{equation}
If it is solvable and the nonzero Smith invariant factors of $M$ are
$d_1,\ldots,d_r$, then its solution set has dimension $n-r$ and has
$d_1\cdots d_r$ connected components.
\end{lemma}
\begin{proof}
Choose unimodular matrices $U\in\operatorname{GL}(m,\Z)$ and
$V\in\operatorname{GL}(n,\Z)$ such that
\[
 UMV=\operatorname{diag}(d_1,\ldots,d_r,0,\ldots,0),
 \qquad d_i>0.
\]
The induced coordinate changes are automorphisms of the source and target
tori.  In these coordinates, the first $r$ equations are
$d_ix_i=(Uy)_i$ in $\R/\Z$ and are always solvable, while the remaining
equations require $(Uy)_i=0$ in $\R/\Z$ for $i>r$.  These conditions are
precisely \eqref{eq:left-kernel-test}, because the last $m-r$ rows of $U$
form an integral basis of $\ker(M^T)\cap\Z^m$.  Each of the first $r$
equations has $d_i$ solutions, and the remaining $n-r$ source coordinates
are free.  This gives the asserted dimension and number of components.
\end{proof}

Every $g\in G$ is uniquely $F^a\iota_1^b\iota_2^c\kappa^d$.  The complete
list of $g\tau$ with fixed points is Table~\ref{tab:coset-fixed}; signs are
in the coordinate order \eqref{eq:seven-coordinates}.  The $s$-coordinate
of $gT$ has sign $-$, so adjoining it neither changes the displayed fixed
dimension nor introduces an additional free coordinate.

\pagebreak[4]
\begin{longtable}{@{}c>{\raggedright\arraybackslash}p{.40\textwidth}c c@{}}
\caption{Fixed elements in $G\tau$.}\label{tab:coset-fixed}\\
\toprule $g\tau$&translation $v$&sign tuple&$\dim\Fix$\\\midrule\endfirsthead
\toprule $g\tau$&translation $v$&sign tuple&$\dim\Fix$\\\midrule\endhead
$\tau$&$(\frac14,0,0,0,0,\frac14,0)$&$(-,+,+,+,+,-,-)$&4\\
$\iota_1\tau$&$(\frac14,0,0,0,0,\frac34,0)$&$(-,-,-,-,-,-,-)$&0\\
$\iota_2\tau$&$(\frac14,0,0,0,0,\frac14,\frac12)$&$(-,-,-,-,-,-,-)$&0\\
$\kappa\tau$&$(0,0,0,0,0,\frac14,0)$&$(+,+,-,+,-,-,+)$&4\\
$\iota_1\kappa\tau$&$(0,0,0,0,0,\frac34,0)$&$(+,-,+,-,+,-,+)$&4\\
$\iota_1\iota_2\tau$&$(\frac14,0,0,0,0,\frac34,\frac12)$&$(-,+,+,+,+,-,-)$&4\\
$F^2\iota_1\iota_2\tau$&$(\frac34,0,0,0,0,\frac34,\frac12)$&$(-,-,-,-,-,-,-)$&0\\
$F^2\iota_1\tau$&$(\frac34,0,0,0,0,\frac34,0)$&$(-,+,+,+,+,-,-)$&4\\
$F^2\tau$&$(\frac34,0,0,0,0,\frac14,0)$&$(-,-,-,-,-,-,-)$&0\\
$F^2\iota_2\tau$&$(\frac34,0,0,0,0,\frac14,\frac12)$&$(-,+,+,+,+,-,-)$&4\\
\bottomrule
\end{longtable}

\begin{lemma}\label{lem:coset-exhaustive}
Table~\ref{tab:coset-fixed} is exhaustive.  Fourteen nonidentity elements
of $K$ have a four-dimensional fixed-point set on $\T^8$.  Four have a
zero-dimensional fixed-point set.  Every remaining nonidentity element
acts freely.
\end{lemma}
\begin{proof}
We write $g=F^a\iota_1^b\iota_2^c\kappa^d$.  If $d=0$ and $a$ is odd, the
$x_3$ equation is translation by an odd quarter-period.  If $d=0$ and $a$
is even, all eight choices have fixed points; four have linear part $-I_7$
and four have a four-dimensional fixed space, exactly as listed.  For
$d=1$, solvability of the $\bar t$ equation forces $a=0$, and solvability
of the $y_3$ equation then forces $c=0$; the only fixed elements in this
case are $\kappa\tau$ and $\iota_1\kappa\tau$.  Thus six elements of
$G\tau$ have four-dimensional fixed sets and four elements of $G\tau$ have
zero-dimensional fixed sets.  The eight fixed elements of $G$ have
three-dimensional fixed sets on $\T^7$, hence four-dimensional fixed sets
after adjoining the fixed $s$-coordinate.  Therefore $K$ has fourteen
nonidentity elements with four-dimensional fixed sets and four with
zero-dimensional fixed sets.
\end{proof}

\subsection{The complete isotropy hierarchy}
Each of the four elements of $G\tau$ whose fixed set is zero-dimensional
has linear part $-I_7$; the corresponding element of $GT$ therefore has
linear part $-I_8$ and fixes $2^8=256$ points of $\T^8$.  These four
fixed-point sets are pairwise disjoint.  For the four new elements listed
in the left column below, the right column lists the unique nonidentity
element of $G$ that fixes the same points:
\begin{equation}\label{eq:zero-old-stabilizer}
\begin{array}{c|c}
\iota_1T&\kappa\\ \iota_2T&\iota_2\kappa\\
F^2\iota_1\iota_2T&F^2\iota_2\kappa\\ F^2T&F^2\kappa.
\end{array}
\end{equation}
Proposition~\ref{prop:isotropy} shows that the resulting order-four groups
are the full point stabilizers.  Every point orbit therefore has
$64/4=16$ elements, and the $4\cdot256$ fixed points determine
$1024/16=64$ isolated points in $\T^8/K$.

Four subgroups have two-dimensional common fixed sets:
\begin{align}
A_0&=\langle\kappa,T\rangle,&
A_1&=\langle\iota_2\kappa,\iota_1\iota_2T\rangle,\label{eq:A01}\\
A_2&=\langle\kappa T,F^2\iota_2\kappa\rangle,&
A_3&=\langle F^2\kappa,\iota_1\kappa T\rangle.\label{eq:A23}
\end{align}
Each common fixed set is a disjoint union of $64$ two-tori.  The subgroups
form two $K$-conjugacy classes of size two.

\begin{proposition}\label{prop:isotropy}
For every $p\in\T^8$ with nontrivial stabilizer $K_p$, exactly one of the
following three alternatives holds.  First, $K_p\cong\Z_2$ and acts as
$(I_4,-I_4)$.  Second, $p$ lies on a two-dimensional stratum,
$K_p\cong\Z_2^2$, and the three nonidentity elements act on the
six-dimensional normal space by the nontrivial even sign changes; the
local model is
\begin{equation}\label{eq:SU3-local}\R^2\times\C^3/\Z_2^2.\end{equation}
Third, $p$ is one of the $64$ isolated quotient points,
$K_p\cong\Z_2^2$, and the local model is
\begin{equation}\label{eq:product-local}
(\R^4/\{\pm1\})\times(\R^4/\{\pm1\}).
\end{equation}
There are no larger stabilizers.
\end{proposition}
\begin{proof}
For affine elements $h_\nu(q)=A_\nu q+v_\nu$, a common fixed point is a
solution of the combined congruence
\begin{equation}\label{eq:stacked}
\begin{pmatrix}A_1-I\\\vdots\\A_r-I\end{pmatrix}q
=-\begin{pmatrix}v_1\\\vdots\\v_r\end{pmatrix}\pmod{\Z^{8r}}.
\end{equation}
For each intersecting pair $(h_1,h_2)$, we record the fixed-dimension data
\[
\bigl(\dim\Fix h_1,\dim\Fix h_2,\dim\Fix(h_1h_2);
      \dim(\Fix h_1\cap\Fix h_2)\bigr).
\]
Lemma~\ref{lem:torus-congruence} applies to the block matrix in
\eqref{eq:stacked}.  Its criterion and Smith normal form give
exactly twelve pairs with data $(4,4,4;2)$.  These pairs assemble three at
a time into the four subgroups \eqref{eq:A01}--\eqref{eq:A23}.  For each
subgroup, the Smith normal form has component group of order $64$, so its
common fixed set lifts to $64$ two-tori.  Every other nonempty pairwise
intersection belongs to one of the four order-four point stabilizers in
\eqref{eq:zero-old-stabilizer}.  Repeating the calculation for triples
yields no additional system: there are exactly four systems of type
$(4,4,4)$ and four of type $(4,4,0)$.  Testing the remaining elements of
$K$ on these solution sets shows that none fixes a point of any set, so no
stabilizer is larger.

At a two-dimensional stratum, the normal representation is the direct sum
of the three nontrivial real characters of $\Z_2^2$, each with multiplicity
two.  After choosing a complex coordinate on each character plane, the
three nonidentity elements act by the even sign changes in
\eqref{eq:SU3-local}.  At an isolated point, we denote the unique
nonidentity stabilizer element in $G$ by $a$ and the element of $GT$ with
zero-dimensional fixed set by $b$.  The eigenspace decomposition of $a$
gives coordinates in which
\[
a=(I_4,-I_4),\qquad b=(-I_4,-I_4),\qquad ab=(-I_4,I_4),
\]
so quotienting by $\langle a,b\rangle$ gives the product model
\eqref{eq:product-local}.
\end{proof}

\section{Joyce resolution data}
We use Joyce's notation \cite[Definition 13.2.2]{JoyceBook}.  For a
connected fixed component $C$, we define
\[
 N_K(C)=\{\gamma\in K:\gamma C=C\},\qquad
 C_K(C)=\{\gamma\in K:\gamma|_C=\id_C\}.
\]
We denote the generic exact isotropy along $C$ by $A_C$, its tangent space
by $V_C$, and its orthogonal complement by $W_C=V_C^\perp$.  Because $C$ is a connected component of the
fixed set of $A_C$ and its generic points have no larger stabilizer,
$C_K(C)=A_C$.  Joyce's residual component group is therefore
\begin{equation}\label{eq:component-group}
 B_C=N_K(C)/C_K(C)=N_K(C)/A_C.
\end{equation}
This is different from the subgroup normalizer $N_K(A_C)$: in general
$N_K(C)\subset N_K(A_C)$.  If $C=\{p\}$ is a point, then
$N_K(p)=C_K(p)$ and $B_p=1$.

\subsection{The four-dimensional strata}
We use the coframe
\begin{equation}\label{eq:e-coframe}
(e^0,e^1,e^2,e^3,e^4,e^5,e^6,e^7)
=(ds,dt,dx_1,dy_1,dx_2,dy_2,dx_3,dy_3).
\end{equation}
We write $(e_0,\ldots,e_7)$ for its dual frame.  For an oriented normal
basis $(w_1,\ldots,w_4)$, we define
$\sigma_1=w^{12}+w^{34}$, $\sigma_2=w^{13}-w^{24}$,
$\sigma_3=w^{14}+w^{23}$.  We also define
$p_\pm=(e_2\pm e_3)/\sqrt2$, $q_\pm=(e_4\pm e_5)/\sqrt2$.

\begin{longtable}{@{}c c c >{\raggedright\arraybackslash}p{.27\textwidth}
>{\raggedright\arraybackslash}p{.34\textwidth}@{}}
\caption{Normal data for the six four-stratum conjugacy classes.}
\label{tab:rdata4}\\
\toprule $h$&$|[h]|$&$|Z_K(h)|$&$V_h$&$W_h$ and $\zeta_h$\\\midrule\endfirsthead
\toprule $h$&$|[h]|$&$|Z_K(h)|$&$V_h$&$W_h$ and $\zeta_h$\\\midrule\endhead
$\kappa$&2&32&$(e_0,e_2,e_4,e_6)$&$(e_1,e_7,e_3,e_5)$; $e^{17}+e^{35}$\\
$T$&2&32&$(e_2,e_3,e_4,e_5)$&$(e_0,e_6,e_1,-e_7)$; $e^{06}-e^{17}$\\
$\kappa F$&4&16&$(e_0,p_-,q_-,e_7)$&$(e_1,-e_6,p_+,q_+)$;
$-e^{16}+\frac12(e^2+e^3)\wedge(e^4+e^5)$\\
$\iota_2\kappa$&2&32&$(e_0,e_3,e_5,-e_6)$&$(e_1,-e_7,e_2,e_4)$; $-e^{17}+e^{24}$\\
$\kappa T$&2&32&$(e_1,e_2,e_4,-e_7)$&$(e_0,-e_6,e_3,e_5)$; $-e^{06}+e^{35}$\\
$\iota_1\iota_2T$&2&32&$(e_2,e_3,e_4,e_5)$&$(e_0,e_6,e_1,-e_7)$; $e^{06}-e^{17}$\\
\bottomrule
\end{longtable}

The rows of Table~\ref{tab:rdata4} classify conjugacy classes of
generators, rather than connected fixed components.  We denote by
$\mathcal F_4$ the finite set of connected components with generic exact isotropy
$\Z_2$.  For each $K$-orbit in $\mathcal F_4$, we choose one component.  Its
generator is conjugate to exactly one row of the table, and the displayed
linear data are attached to that component and then transported around its
orbit.  We make the same orbitwise choice for the components of the four
two-dimensional common fixed sets and for the isolated points.  Transporting
these choices assigns resolution data to every connected component of every
stratum.

\begin{lemma}\label{lem:four-rdata}
Each $V_h$ is Cayley-oriented, $\zeta_h$ is self-dual on $W_h$, and the
image of $Z_K(h)$ on $\Lambda^+W_h^*$ is
\begin{equation}\label{eq:centralizer-image}
\{I_3,\diag(1,-1,-1)\}
\end{equation}
in a basis beginning with $\zeta_h$.
\end{lemma}
\begin{proof}
The expanded Cayley form is
\begin{align}\label{eq:Phi-expanded}
\Phi_0={}&e^{0123}+e^{0145}+e^{0167}+e^{0246}-e^{0257}-e^{0347}-e^{0356}\\
&+e^{2345}+e^{2367}+e^{4567}-e^{1247}-e^{1256}-e^{1346}+e^{1357}.\nonumber
\end{align}
For each row of Table~\ref{tab:rdata4}, evaluation of the explicit form
\eqref{eq:Phi-expanded} on the listed oriented basis of $V_h$ is $+1$;
hence that orientation is the Cayley orientation.  Relative to the listed
oriented basis of $W_h$, the form $\zeta_h$ equals $w^{12}+w^{34}$, so it
is self-dual and is the first vector of the standard oriented basis of
$\Lambda^+W_h^*$.  We restrict the five generator matrices to $W_h$ and
write a general element in the form
$q=F^a\iota_1^b\iota_2^c\kappa^dT^e$.  For
$h=\kappa,T,\iota_2\kappa,\kappa T,\iota_1\iota_2T$, the equation
$qh=hq$ is exactly $a\equiv0\pmod2$, giving $32$ elements.  For
$h=\kappa F$ it is exactly
$e=0$ and $d\equiv a\pmod2$, giving $16$ elements.  Thus the centralizer
orders are $32$ and $16$, respectively, and orbit--stabilizer gives the
class sizes displayed in Table~\ref{tab:rdata4}.  The induced action of
each centralizer on $(\sigma_1,\sigma_2,\sigma_3)$ has image
$\{I_3,\diag(1,-1,-1)\}$.  The nontrivial image is realized by $\kappa$
for $h=T$, by $\iota_1$ for $h=\kappa F$, and by the corresponding
directly computed centralizer element in each remaining row.
\end{proof}

We fix numbers $L>2\ell>0$.  For a four-dimensional component $C$ with
pointwise stabilizer generated by $h_C$, we define
\begin{equation}\label{eq:four-scale}
s(C)=\begin{cases}L,&h_C\in G,\\ \ell,&h_C\in GT.\end{cases}
\end{equation}
We equip the radial Eguchi--Hanson resolution
$Y_C=\EH_{s(C)}\to W_C/\{\pm1\}$ with the normalized K\"ahler form whose
ALE limit lies on the ray $\R_{>0}\zeta_C$ and whose integral over the
exceptional sphere is $s(C)$.  To specify the full product $\SpinSeven$
structure, we must also fix the remaining two members of its hyperk\"ahler
triple.  We write $(v_1,\ldots,v_4)$ for the ordered frame of $V_C$ in
Table~\ref{tab:rdata4}, and we define
\[
 \alpha_1=v^{12}+v^{34},\qquad
 \alpha_2=v^{13}-v^{24},\qquad
 \alpha_3=v^{14}+v^{23}.
\]
We choose the oriented self-dual basis
$(\sigma_1,\sigma_2,\sigma_3)$ of $W_C^*$ beginning with
$\sigma_1=\zeta_C$ for which the flat splitting of
\eqref{eq:Phi-expanded} is
\begin{equation}\label{eq:four-product-flat}
 \Phi_0=v^{1234}+\alpha_1\wedge\sigma_1
       +\alpha_2\wedge\sigma_2-\alpha_3\wedge\sigma_3
       +\tfrac12\sigma_1^2.
\end{equation}
We denote by $(\omega_{C,1},\omega_{C,2},\omega_{C,3})$ the parallel
hyperk\"ahler triple on $\EH_{s(C)}$ that is asymptotic to
$(\sigma_1,\sigma_2,\sigma_3)$ and whose class
$[\omega_{C,1}]$ evaluates to $s(C)$ on the exceptional sphere.  The
complete product model used below is therefore
\begin{equation}\label{eq:four-product-EH}
 \Phi_C=v^{1234}+\alpha_1\wedge\omega_{C,1}
       +\alpha_2\wedge\omega_{C,2}-\alpha_3\wedge\omega_{C,3}
       +\tfrac12\omega_{C,1}^2,
 \qquad g_C=g_{V_C}+g_{\EH_{s(C)}}.
\end{equation}
Equation \eqref{eq:four-product-EH} is Joyce's product formula from
\cite[Proposition 13.1.1]{JoyceBook} and specifies the complete metric and
Cayley form at the four-dimensional stratum.
Because $G\triangleleft K$, conjugation preserves the alternatives in
\eqref{eq:four-scale}.  Thus $s(C)=s(C')$ whenever $\gamma C=C'$.  If
$\gamma C=C'$, we denote by $\rho_+(D\gamma|_{W_C})$ the transport on
self-dual two-forms induced by the oriented isometry
$D\gamma:W_C\to W_{C'}$, and we define
\begin{equation}\label{eq:epsilon}
\zeta_{C'}=\epsilon_{\gamma,C}\rho_+(D\gamma|_{W_C})\zeta_C,
\qquad \epsilon_{\gamma,C}\in\{\pm1\}.
\end{equation}
If $\epsilon_{\gamma,C}=+1$, we lift $D\gamma|_{W_C}$ by the radial
holomorphic Eguchi--Hanson isometry.  If $\epsilon_{\gamma,C}=-1$, we use
the radial antiholomorphic lift.  Definition \eqref{eq:epsilon} gives the
cocycle identity
\begin{equation}\label{eq:epsilon-cocycle}
\epsilon_{\gamma\delta,C}=\epsilon_{\gamma,\delta C}\epsilon_{\delta,C}.
\end{equation}
Lemma~\ref{lem:four-rdata} implies $\epsilon_{\gamma,C}=1$ for every
$\gamma\in N_K(C)$, so transport from an orbit representative is
independent of the transporting element.  Each of the two possible lifts
preserves the Eguchi--Hanson metric and induces on its hyperk\"ahler triple
the same orthogonal transformation that $D\gamma$ induces on
$\Lambda^+W_C^*$.  Because $D\gamma$ preserves the flat identity
\eqref{eq:four-product-flat}, its action on
$(\alpha_1,\alpha_2,\alpha_3)$ is paired with the normal action so that
every term in \eqref{eq:four-product-EH} is fixed.  Therefore the
transported four-stratum metric and Cayley form are exactly
$K$-equivariant.

\subsection{The two- and zero-dimensional strata}
For $A_0=\langle\kappa,T\rangle$, $V_{A_0}=\langle e_2,e_4\rangle$ and a
unitary normal coframe is
\begin{equation}\label{eq:SU3-coframe0}
\xi_1=e^0+\ii e^6,\quad \xi_2=e^3-\ii e^5,\quad \xi_3=e^1-\ii e^7.
\end{equation}
Relative to $(\xi_1,\xi_2,\xi_3)$, the elements $\kappa$ and $T$ act by
the sign triples $(1,-1,-1)$ and $(-1,1,-1)$, respectively.  For $A_1$,
whose common plane is $\langle e_3,e_5\rangle$, we use the unitary normal
coframe
\begin{equation}\label{eq:SU3-coframe1}
\eta_1=e^0-\ii e^6,\quad \eta_2=e^2-\ii e^4,\quad \eta_3=e^1+\ii e^7.
\end{equation}
The subgroups $A_2$ and $A_3$ are $K$-conjugate to $A_1$ and $A_0$,
respectively.  Hence every two-stratum normal representation is conjugate
in $SU(3)$ to the diagonal subgroup
\begin{equation}\label{eq:standard-A}
A=\{(1,1,1),(1,-1,-1),(-1,1,-1),(-1,-1,1)\}\subset SU(3).
\end{equation}
The corresponding flat Calabi--Yau forms are
\begin{align}
\omega_{A_0}&=e^{06}-e^{35}-e^{17},&
\Omega_{A_0}&=\xi_1\wedge\xi_2\wedge\xi_3,\label{eq:CY-A0}\\
\omega_{A_1}&=-e^{06}-e^{24}+e^{17},&
\Omega_{A_1}&=\eta_1\wedge\eta_2\wedge\eta_3.\label{eq:CY-A1}
\end{align}
We choose the oriented tangent coframe $(u^1,u^2)=(-e^2,-e^4)$ for the
$A_0$-stratum and $(u^1,u^2)=(e^3,e^5)$ for the $A_1$-stratum.  Expanding
the right-hand side of the following identity with
\eqref{eq:CY-A0}--\eqref{eq:CY-A1} reproduces every term of
\eqref{eq:Phi-expanded}; therefore, for $r=0,1$,
\begin{equation}\label{eq:flat-SU3-product}
\Phi_0=u^{12}\wedge\omega_{A_r}
 +\operatorname{Re}\bigl((u^1+\ii u^2)\wedge\Omega_{A_r}\bigr)
 +\tfrac12\omega_{A_r}^2.
\end{equation}
Suppose that a component fixed by $A_r$ lies in a four-dimensional
component whose generic stabilizer is generated by $h\in A_r$.  We denote
the complex normal coordinate fixed by $h$ by $\lambda_h$, and we denote
the direct sum of the other two complex coordinate lines by $Q_h$.  If
$Z_h$ is the unique coordinate $(1,0)$-vector satisfying
$d\lambda_h(Z_h)=1$, we define the complete ordered flat transverse triple
by
\begin{equation}\label{eq:incidence-triple}
 \varpi_{h,1}=\omega_{A_r}|_{Q_h},\qquad
 \varpi_{h,2}+\ii\varpi_{h,3}
   =(\iota_{Z_h}\Omega_{A_r})|_{Q_h}.
\end{equation}
Equivalently, $\varpi_{h,1}$ is the restriction of $\omega_{A_r}$ to
$Q_h$, while $\varpi_{h,2}$ and $\varpi_{h,3}$ are the real and imaginary
parts of the two-form obtained by contracting $\Omega_{A_r}$ with $Z_h$.
With
$\zeta_{\iota_1\kappa T}=-e^{06}-e^{24}$ (using the oriented normal basis
$(e_0,-e_6,e_2,-e_4)$), the six restrictions are
\begin{equation}\label{eq:incidence-signs}
\begin{array}{c|c|c|c}
 &\multicolumn{3}{c}{h:\ \text{fixed complex line},\ \omega_{\perp,h}}\\
\hline
A_0&\kappa:\xi_1,-\zeta_\kappa&T:\xi_2,+\zeta_T&
\kappa T:\xi_3,-\zeta_{\kappa T}\\
A_1&\iota_2\kappa:\eta_1,-\zeta_{\iota_2\kappa}&
\iota_1\iota_2T:\eta_2,-\zeta_{\iota_1\iota_2T}&
\iota_1\kappa T:\eta_3,+\zeta_{\iota_1\kappa T}.
\end{array}
\end{equation}
In each row the first involution belongs to $G$ and its transverse
Eguchi--Hanson factor has scale $L$; the other two belong to $GT$ and have
scale $\ell$.  For every incidence in \eqref{eq:incidence-signs}, the line
$\R\varpi_{h,1}$ equals the distinguished line $\R\zeta_h$ of the incident
four-stratum, and the table records the orientation sign.  When the sign
is positive, we identify the transverse resolution with the incident
Eguchi--Hanson space by the radial holomorphic lift.  When the sign is
negative, we use the radial antiholomorphic lift and choose the phase of
its holomorphic symplectic form so that the pullback of
$(\omega_{C,1},\omega_{C,2},\omega_{C,3})$ equals
$(\varpi_{h,1},\varpi_{h,2},\varpi_{h,3})$.  Both triples are parallel and
have the same ALE limit under this identification, so their difference is
a parallel tensor that tends to zero and therefore vanishes identically.
The incidences for $A_3=FA_0F^{-1}$ and $A_2=FA_1F^{-1}$ are defined by
this $F$-transport, and their orientation signs are determined by
\eqref{eq:epsilon-cocycle}.

In each of the coframes \eqref{eq:SU3-coframe0} and
\eqref{eq:SU3-coframe1}, the unique nonidentity element of $A_r\cap G$
acts by $(1,-1,-1)$ and fixes the first complex coordinate line.  We define
$m_{ij}=(e_i+e_j)/2$ and
\begin{equation}\label{eq:toric-lattices}
N_a=\Z^3+\Z m_{23},\qquad
N=N_a+\Z m_{13}
  =\Z^3+\Z m_{12}+\Z m_{13}.
\end{equation}
We first quotient by the old subgroup
$A_r\cap G=\langle(1,-1,-1)\rangle$.  Its crepant subdivision is taken in
$N_a$ and has the two maximal cones
\[
\langle e_1,e_2,m_{23}\rangle,\qquad
\langle e_1,m_{23},e_3\rangle.
\]
Passing to the residual involution replaces $N_a$ by the quotient lattice
$N$.  The residual fixed curves correspond to the two faces
$\langle e_1,e_2\rangle$ and $\langle e_1,e_3\rangle$; star subdivision of
these faces introduces the primitive rays $m_{12}$ and $m_{13}$ and gives
the staged fan
\begin{equation}\label{eq:toric-cones}
\begin{split}
\Sigma_{\rm st}^{\max}=\{&
\langle e_1,m_{12},m_{23}\rangle,\
\langle e_2,m_{12},m_{23}\rangle,\\
&\langle e_1,m_{13},m_{23}\rangle,\
\langle e_3,m_{13},m_{23}\rangle\}.
\end{split}
\end{equation}
For each of the four maximal cones, the determinant of its primitive ray
generators in the standard lattice is $\pm1/4$, which equals the covolume
of $N$.  Those generators therefore form a basis of $N$, and the fan is
smooth.  Every primitive ray generator lies on the height-one hyperplane,
so the subdivision is crepant.  Because the subdivision has the same
support $\R_{\geq0}^3$ as the affine quotient fan, the associated toric
morphism is proper and birational.  Hence
\begin{equation}\label{eq:Y-A}
Y_A=Y_{\Sigma_{\rm st}}\longrightarrow\C^3/A
\end{equation}
is smooth and crepant.  By \cite[Proposition 9.1.4]{JoyceBook}, every
crepant resolution of $\C^m/H$ is a local-product resolution in the sense
required by Joyce's $R$-data definition.

We denote generic exceptional $\CP^1$ fibres transverse to the first,
second, and third coordinate axes by $C_{23},C_{13},C_{12}$,
respectively.  Consider the torus-invariant real divisor
\[
 D=-\frac L2D_{m_{23}}
   -\frac\ell2\bigl(D_{m_{12}}+D_{m_{13}}\bigr).
\]
In Fulton's convention, the support function of
$D=\sum_\rho a_\rho D_\rho$ satisfies
$\psi_D(v_\rho)=-a_\rho$.  Thus the support function associated with $D$
is defined on the primitive rays of $\Sigma_{\rm st}$ by
\begin{equation}\label{eq:support-values}
\psi_D(e_i)=0,\qquad \psi_D(m_{23})=\tfrac L2,\qquad
\psi_D(m_{12})=\psi_D(m_{13})=\tfrac\ell2.
\end{equation}
The compact walls
$\langle m_{12},m_{23}\rangle$,
$\langle m_{13},m_{23}\rangle$, and
$\langle e_1,m_{23}\rangle$ are exactly the internal two-cones of the fan
and have wall degrees $\ell$, $\ell$, and $(L-2\ell)/2$, respectively.
For the third wall, this computation uses
$m_{12}+m_{13}=e_1+m_{23}$.  Since $L>2\ell>0$, all three wall degrees are
positive, and $\psi_D$ is strictly upper convex in Fulton's convention.
Restriction to the three
transverse $A_1$ fans, whose primitive relations are
$e_i+e_j=2m_{ij}$, gives the QALE K\"ahler class periods
\begin{equation}\label{eq:toric-periods}
\int_{C_{23}}\omega_A=L,\qquad
\int_{C_{13}}\omega_A=\int_{C_{12}}\omega_A=\ell.
\end{equation}
We normalize the compact torus as $N_{\R}/N$, so every primitive angular
coordinate has period one.  In a convention in which primitive angles
have period $2\pi$, the support values must instead be divided by $2\pi$;
with either convention the geometric periods are those in
\eqref{eq:toric-periods}.  The toric support-function theorem
\cite[\S3.4]{Fulton} says
that strict upper convexity is the toric ampleness criterion.  In the
present noncompact situation this is first a relative statement for
$Y_A\to\C^3/A$, because the fan has the support of the single affine cone.
Since the base is affine, a relatively ample rational divisor is ample on
$Y_A$; openness of the ample cone and linearity extend this conclusion to
the real divisor $D$.  Its real first Chern class is therefore a K\"ahler
class, and its integrals over invariant compact curves are the wall
pairings above.  The pair $(A,\C^3/A)$ is the one treated
in \cite[Example 9.3.6]{JoyceBook}, which proves that each of its four
crepant resolutions admits a QALE K\"ahler metric.  The toric variety
$Y_A$ is one of these resolutions and is local-product by Proposition
9.1.4.  Joyce's \cite[Theorem 9.4.7]{JoyceBook} therefore applies to the
specific K\"ahler class defined by $D$, equivalently by
\eqref{eq:support-values}, and supplies
a QALE K\"ahler representative of that class.

We apply Joyce's Theorem 9.3.3 to the K\"ahler class just constructed and
denote its unique Ricci-flat QALE representative by $g_A$, with K\"ahler
form $\omega_A$.  We normalize the parallel holomorphic volume form
$\Omega_A$ so that its limit is \eqref{eq:CY-A0} or \eqref{eq:CY-A1},
according to the two-stratum orbit.  The resulting Calabi--Yau threefold
$(Y_A,g_A,\omega_A,\Omega_A)$ has one singular-axis end asymptotic to
$\C\times\EH_L$ and two ends asymptotic to $\C\times\EH_\ell$, with the
hyperk\"ahler identifications in \eqref{eq:incidence-signs}.

The existence of $\Omega_A$ follows from the crepant toric construction.
The invariant flat volume form on the regular locus of $\C^3/A$ pulls back
to the inverse image of that locus in $Y_A$.  The height-one condition
implies that its divisor along every toric exceptional divisor is zero, so
it extends to a nowhere-vanishing holomorphic three-form on all of $Y_A$.
Joyce's \cite[Example 9.3.6 and Theorem 9.3.4]{JoyceBook} identifies the
holonomy of $g_A$ with $SU(3)$; hence the extended volume form is parallel
and unique up to a nonzero constant, whose modulus and phase are fixed by
the prescribed flat-end limit.  Proposition 13.1.2 and Theorem 13.1.10(a)
of \cite{JoyceBook} associate to these Calabi--Yau data the required
torsion-free product $\SpinSeven$ structure on $\R^2\times Y_A$.

At an isolated stratum, we denote the nonidentity old stabilizer by
$a\in G$ and the all-minus stabilizer element by $b\in GT$.  We choose the
eigenspace splitting $\R^8=U\oplus U'$ for which
$a=(I_U,-I_{U'})$, $b=(-I_U,-I_{U'})$, and
$ab=(-I_U,I_{U'})$, and we use the product resolution
\begin{equation}\label{eq:EH-product}
\EH_\ell(U)\times\EH_L(U')\longrightarrow
(U/\{\pm1\})\times(U'/\{\pm1\}).
\end{equation}
We choose quaternionic structures on $U$ and $U'$ that are compatible with
$\Phi_0$.  Under these identifications, the stabilizer is the sign subgroup
of $Sp(1)\times Sp(1)\subset Sp(2)$.  We choose parallel triples
$(\omega_1,\omega_2,\omega_3)$ and
$(\omega'_1,\omega'_2,\omega'_3)$ on the two Eguchi--Hanson factors and
normalize their flat limits so that the following expression tends to
\eqref{eq:Phi-expanded}:
\begin{equation}\label{eq:point-Cayley}
\begin{split}
 \Phi_p
 &=\tfrac12(\omega_1+\omega'_1)^2
   +\operatorname{Re}\bigl((\omega_2+\ii\omega_3)
                      \wedge(\omega'_2+\ii\omega'_3)\bigr)\\
 &=\operatorname{vol}_{\EH_\ell}+\operatorname{vol}_{\EH_L}
   +\omega_1\wedge\omega'_1+
    \omega_2\wedge\omega'_2-\omega_3\wedge\omega'_3.
\end{split}
\end{equation}
The Riemannian metric associated with $\Phi_p$ is the product metric
\begin{equation}\label{eq:point-metric}
 g_p=g_{\EH_\ell(U)}+g_{\EH_L(U')}.
\end{equation}
Equation \eqref{eq:point-Cayley} is Proposition 13.1.4 applied to the
product Calabi--Yau four-fold,
whose K\"ahler form is $\omega_1+\omega'_1$ and whose holomorphic volume
form is $(\omega_2+\ii\omega_3)\wedge
(\omega'_2+\ii\omega'_3)$.  Hence $\Phi_p$ is closed and torsion-free.
It remains to verify the QALE estimates for the unequal-scale product.
Write $r_U=1+|u|$ and $r_{U'}=1+|u'|$.  On the end modelled on
$U\times\EH_L(U')$, the ALE chart replaces the $\EH_\ell(U)$ factor by
its flat cone and gives, for every $k\geq0$,
\[
 \nabla^k(\Phi_p-\Phi_{U\times\EH_L})=O(r_U^{-4-k});
\]
on the end modelled on $\EH_\ell(U)\times U'$, one similarly has
\[
 \nabla^k(\Phi_p-\Phi_{\EH_\ell\times U'})=O(r_{U'}^{-4-k}).
\]
We compare these estimates with the weights in
\cite[Definition 13.1.8]{JoyceBook}.  In the first chart the index $i$
has fixed space $V_i=U$.  Let $j$ be the index of the complementary fixed
space $V_j=U'$.  Then $W_j=U$ has real dimension four, so Joyce's exponent
is $d_j=2-\dim_{\R}W_j=-2$.  The subset $V_j/A_i$ of
$V_i\times W_i/A_i$ is $\{0\}\times U'/\{\pm1\}$; hence
$\mu_{i,j}\asymp r_U$ on the QALE region.  The minimum defining $\nu_i$
also contains $\mu_{i,i}$, which measures the transverse radius and can
be smaller than $r_U$.  What is needed is the one-sided estimate
$1\leq\nu_i\leq1+\mu_{i,j}\leq C r_U$.  Since the exponents are negative,
this gives
\[
 r_U^{-4-k}\leq C
 \mu_{i,j}^{d_j}\nu_i^{-2-k}.
\]
Thus the first ALE error is $O$ of the weight required by Joyce's
definition, even when the required weight is larger than
$r_U^{-4-k}$.  The second chart is identical after interchanging $U$ and
$U'$.

In the ambient chart, the two four-plane indices satisfy
\[
 \mu_{\mathbf0,U}\asymp r_{U'},\qquad
 \mu_{\mathbf0,U'}\asymp r_U,\qquad
 \nu_{\mathbf0}\asymp1+\min(|u|,|u'|),
\]
and both exponents are $-2$.  If, for example, $r_U\geq r_{U'}$, then
every derivative of the error contributed by the $U$-factor, every
derivative of the error contributed by the $U'$-factor, and every mixed
product of the two errors is bounded by a constant multiple of
$r_{U'}^{-4-k}$.  This is the summand
$O(\mu_{\mathbf0,U}^{-2}\nu_{\mathbf0}^{-2-k})$; the case
$r_{U'}\geq r_U$ is symmetric.  Thus, for every $k\geq0$, the full error
is bounded by
\[
 \sum_{q:\,\mathbf0\not\succeq q}
 O\!\left(\mu_{\mathbf0,q}^{d_q}
              \nu_{\mathbf0}^{-2-k}\right).
\]
The remaining, deeper-stratum summands on the right are harmless because
the asserted estimate is an upper bound by the whole sum.  This verifies
Definition 13.1.8 in every chart for
$\EH_\ell\times\EH_L$.

\begin{proposition}\label{prop:R-data}
The four-stratum Eguchi--Hanson models, the two-stratum toric models, the
isolated-point product models, and the transition and transport maps
specified above and in the proof below constitute $R$-data for
$\T^8/K$ in the sense of
\cite[Definition 13.2.2]{JoyceBook}.
\end{proposition}
\begin{proof}
We verify the three compatibility requirements of Joyce's definition.
Condition (i) is verified separately in each stratum dimension.  At a
four-stratum we use the product $\R^4\times\EH_{s(C)}$ of
\cite[Proposition 13.1.1]{JoyceBook}.  At a two-stratum we use
$\R^2\times Y_A$ and \cite[Proposition 13.1.2]{JoyceBook}.  At an isolated
point we use the product \eqref{eq:EH-product} and
\cite[Proposition 13.1.4]{JoyceBook}.  The estimates following
\eqref{eq:EH-product} prove directly that the isolated-point model is
QALE, while \cite[Theorem 13.1.10(a)]{JoyceBook} proves the QALE assertion
for $\R^2\times Y_A$.

For condition (ii), the face $\langle e_2,e_3\rangle$ transverse to the
distinguished first axis is subdivided by $m_{23}$ and has period $L$; the
other two coordinate faces are subdivided by $m_{13},m_{12}$ and have
period $\ell$.  Each transverse fan is the minimal $A_1$ fan, so the three
models are respectively $\EH_L,\EH_\ell,\EH_\ell$.
By \cite[Proposition 9.1.4]{JoyceBook}, $Y_A$ is a local-product
resolution.  For each fixed coordinate axis, we choose the map required
by \cite[Definition 9.1.2]{JoyceBook} to be the toric product map obtained
from the star of the corresponding face,
\begin{equation}\label{eq:axis-product-map}
 \psi_{A,h}:
 \mathcal U_{A,h}
 \longrightarrow Y_A.
\end{equation}
Here $h$ ranges over the three nonidentity elements of $A$.  We set
$s_h=L$ when $h\in A\cap G$, and we set $s_h=\ell$ for the other two
elements of $A$.  The set $\mathcal U_{A,h}$ is the complement of the
deeper-stratum set in the product of the coordinate axis with the
transverse toric $A_1$-resolution.  The star construction identifies it
with an open neighborhood of the corresponding end of $Y_A$, and
$\psi_{A,h}$ commutes with the blowdowns.  Proposition 9.1.5 says that
the induced transverse resolution is itself local-product.  We use these
maps, with the holomorphic or
antiholomorphic transverse identification specified after
\eqref{eq:incidence-signs}.  Formula \eqref{eq:incidence-triple} then
identifies the entire ordered hyperk\"ahler triple, and substitution into
\eqref{eq:four-product-EH} and \eqref{eq:SU3-product-Cayley} identifies the
complete Cayley forms.  At an isolated point, the transition to the
four-stratum modelled on $U\times\EH_L(U')$ is the product of the inverse
ALE chart
\[
 (U/\{\pm1\})\setminus B_R\longrightarrow\EH_\ell(U)
\]
with the identity on $\EH_L(U')$.  The transition to the end modelled on
$\EH_\ell(U)\times U'$ uses the inverse ALE chart
$(U'/\{\pm1\})\setminus B_R\to\EH_L(U')$ in the second factor and the
identity in the first factor.  The normal factors have scales $\ell$ and
$L$, matching the new and old incident four-strata, respectively.

We next identify the lower-dimensional metrics occurring in the QALE
asymptotics.  In a QALE chart
along one of the three singular axes, Definition 9.2.1 supplies a
lower-dimensional QALE K\"ahler metric $g_{A,h}$ on the transverse minimal
$A_1$ resolution; Proposition 9.2.4 says that this asymptotic metric is
itself QALE.  Because $g_A$ is Ricci-flat and the QALE convergence holds
with two derivatives, the Ricci tensors in the rescaled charts converge
to $\operatorname{Ric}(g_{A,h})$; hence $g_{A,h}$ is Ricci-flat.
Restricting \eqref{eq:toric-periods} to the transverse exceptional sphere
gives area $L$ or $\ell$.  As
$H^2(\EH;\R)=\R\,\operatorname{PD}[E]$, where $E$ is the exceptional
sphere, this single period determines the complete K\"ahler class, and the
uniqueness part of Theorem 9.3.3 gives
\[
 g_{A,h}=g_{\EH_L}\quad\hbox{or}\quad g_{A,h}=g_{\EH_\ell}
\]
with the ALE normalization already chosen at the incident four-stratum.
Moreover, the normalized crepant volume form in the local-product chart is
the wedge of the coordinate-axis one-form with the Eguchi--Hanson
holomorphic two-form.  Substitution of this splitting into Joyce's formula
in Proposition 13.1.2 proves the complete all-derivative $\SpinSeven$
asymptotic required by Definition 13.1.8.  For the isolated-point model,
the same all-derivative conclusion follows from the explicit product
estimates above.

In Joyce's partial order, we denote the ambient nonsingular index by
$\mathbf 0$.  Every singular index satisfies $i\succ\mathbf 0$.  The
remaining relations are exactly the following: each two-stratum index lies
above the three incident four-stratum indices in
\eqref{eq:incidence-signs}, and each isolated-point index lies above the two
four-stratum indices associated with $U$ and $U'$.  If a chain contains a
two-stratum, the corresponding transition map is the composition through
the appropriate axis chart \eqref{eq:axis-product-map}.
Proposition~\ref{prop:isotropy} shows that no other stabilizer, and hence no
other comparable pair or chain, occurs.

For condition (iii), we must prove exact equivariance of both the metric
and the Cayley form.  Lemma~\ref{lem:four-rdata} together with
\eqref{eq:epsilon} establishes this equivariance at every four-stratum.
For the representatives $A_0,A_1$, direct multiplication in the normal
form $F^a\iota_1^b\iota_2^c\kappa^dT^e$ shows that membership in either
subgroup normalizer is equivalent to $a\equiv0\pmod2$, which gives
\begin{equation}\label{eq:two-normalizers}
N_K(A_0)=N_K(A_1)
=\{F^a\iota_1^b\iota_2^c\kappa^dT^e:a\equiv0\pmod2\}.
\end{equation}
Thus each subgroup normalizer has order $32$ and is generated by
$F^2,\iota_1,\iota_2,\kappa,T$.  Their complex signs in the coframes
\eqref{eq:SU3-coframe0}--\eqref{eq:SU3-coframe1} are
\begin{equation}\label{eq:two-normalizer-signs}
\begin{array}{c|ccc|ccc}
&\multicolumn{3}{c|}{(\xi_1,\xi_2,\xi_3)}
&\multicolumn{3}{c}{(\eta_1,\eta_2,\eta_3)}\\
\hline
F^2,\iota_1,\iota_2&+&-&+&+&-&+\\
\kappa&+&-&-&+&+&-\\
T&-&+&-&-&+&-
\end{array}
\end{equation}
The derivative of each listed generator is diagonal and complex-linear in
the indicated coframe.  It therefore acts trivially on the fan lattice,
fixes each of the six ray labels in \eqref{eq:toric-cones}, and preserves
every support value in \eqref{eq:support-values}.  Since the normalizer of
a connected component is a subgroup of the subgroup normalizer, these
properties also hold for every element of its component group
$B_{C_A}=N_K(C_A)/A$.  A diagonal complex sign map is an automorphism of
the dense algebraic torus and fixes every cone of
$\Sigma_{\rm st}$; it therefore has a unique holomorphic toric lift
$\widehat\gamma$ to $Y_A$ agreeing with it on that dense torus.

The toric lift $\widehat\gamma$ preserves the K\"ahler class specified by
\eqref{eq:toric-periods}.  Hence $\widehat\gamma^*\omega_A$ is a
Ricci-flat QALE K\"ahler form in the same class and with the same normalized
flat limit as $\omega_A$; uniqueness in Joyce's Theorem 9.3.3 gives
\begin{equation}\label{eq:normalizer-CY}
\widehat\gamma^*\omega_A=\omega_A,\qquad
\widehat\gamma^*\Omega_A=\delta_\gamma\Omega_A,\qquad
\delta_\gamma\in\{\pm1\},
\end{equation}
where $\delta_\gamma$ is the product of the three complex signs in
\eqref{eq:two-normalizer-signs}.  On the flat end, the diagonal complex
signs multiply the holomorphic volume form by $\delta_\gamma$.  Since both
$\widehat\gamma^*\Omega_A$ and $\delta_\gamma\Omega_A$ are parallel
holomorphic volume forms with that same limit, they agree everywhere.

We denote the complex one-form on the tangent two-plane chosen above by
$\nu=u^1+\ii u^2$.  With this choice, the product Cayley form is
\begin{equation}\label{eq:SU3-product-Cayley}
\Phi_A=\tfrac12\omega_A^2+u^{12}\wedge\omega_A
       +\operatorname{Re}(\nu\wedge\Omega_A).
\end{equation}
For each normalizer element $\gamma$, its original affine derivative on
$V_{A_r}\oplus W_{A_r}$ preserves $\Phi_0$.  Together with
$\gamma^*\Omega_{A_r}=\delta_\gamma\Omega_{A_r}$, the flat product identity
forces its tangent action to satisfy
\begin{equation}\label{eq:tangent-phase}
\gamma^*\nu=\delta_\gamma\nu.
\end{equation}
We denote the resulting product lift on $V_{A_r}\times Y_A$ by
\[
 \widetilde\gamma
 =(D\gamma|_{V_{A_r}})\times\widehat\gamma.
\]
Substitution of \eqref{eq:normalizer-CY} and
\eqref{eq:tangent-phase} into \eqref{eq:SU3-product-Cayley} gives
$\widetilde\gamma^*\Phi_A=\Phi_A$ and
$\widetilde\gamma^*((u^1)^2+(u^2)^2+g_A)
=(u^1)^2+(u^2)^2+g_A$ exactly.
At a point, the two involutions with four-dimensional fixed space act as
$(-I_U,I_{U'})$ and $(I_U,-I_{U'})$.  We use the product
$\EH_\ell(U)\times\EH_L(U')$ from \eqref{eq:EH-product}.  Because $B_p=1$,
condition (iii) imposes no nontrivial residual automorphism on the model
attached to a fixed point component; it remains to define transports
between distinct points in the same $K$-orbit.  If
$p'=\gamma p$, normality of $G$ preserves the unique old subgroup
$A_p\cap G$; hence $D\gamma$ carries the labelled character splitting
$U_p\oplus U'_p$ to $U_{p'}\oplus U'_{p'}$ without interchanging the
unequal-scale factors.  We denote the radial lifts of the two linear
quotient maps by
\[
 \widehat{\gamma}_U:\EH_\ell(U_p)\longrightarrow\EH_\ell(U_{p'}),
 \qquad
 \widehat{\gamma}_{U'}:\EH_L(U'_p)\longrightarrow\EH_L(U'_{p'}).
\]
We choose each lift
so that its action on the Eguchi--Hanson hyperk\"ahler triple is the
orthogonal transformation induced by the corresponding restriction of
$D\gamma$ on flat self-dual two-forms; this determines whether the lift is
holomorphic or antiholomorphic.  We define
\begin{equation}\label{eq:point-transport}
 \chi_{\gamma,p}=\widehat{\gamma}_U\times\widehat{\gamma}_{U'}.
\end{equation}
Under these lifts, every term in the bilinear formula
\eqref{eq:point-Cayley} transforms exactly as its flat-limit counterpart.
Since $(D\gamma)^*\Phi_0=\Phi_0$, it follows that
$\chi_{\gamma,p}^*\Phi_{p'}=\Phi_p$, and the radial isometries also give
$\chi_{\gamma,p}^*g_{p'}=g_p$.

For each two-stratum orbit, we define the data on a translate
$C'=\gamma C_A$ by transporting the data on the chosen representative
through $\gamma$.  If two elements give the same translate, their quotient
lies in $N_K(C_A)$, and the proved $B_{C_A}$-invariance shows that the
transported datum is independent of the choice.  We define the
isolated-point data orbitwise by the same procedure.

Every transport map constructed above agrees, away from the exceptional
locus, with the lift of the original affine action.  Therefore, for every
component $C$ and every $\gamma,\delta\in K$, the two sides of
\begin{equation}\label{eq:chi-cocycle}
 \chi_{\gamma\delta,C}
 =\chi_{\gamma,\delta C}\circ\chi_{\delta,C}
\end{equation}
agree on the dense open set on which the blowdown is an isomorphism.  Since
both sides are continuous maps, equality on that dense set implies
equality everywhere.  Thus the transports satisfy Joyce's cocycle
\eqref{eq:chi-cocycle}.  Conditions (i)--(iii), and hence all the
requirements of \cite[Definition 13.2.2]{JoyceBook}, are satisfied.
\end{proof}

\subsection{The global Joyce atlas and analytic gluing}
We now specialize Joyce's Definition 13.2.5 to the full nested isotropy
stratification of $\T^8/K$.  We use $\mathcal I$ to index the distinct connected
submanifolds that occur as components of $\Fix(A)$ for subgroups
$A\leq K$, together with the ambient component
$F_{\mathbf 0}=\T^8$.  The bold index $\mathbf 0$ is distinct from the
two-stratum subgroup $A_0$ in \eqref{eq:A01}.  For each index, we define
\[
 A_i=C_K(F_i),\qquad V_i=T F_i,\qquad W_i=V_i^\perp,
 \qquad B_i=N_K(F_i)/A_i.
\]
The partial order $i\succeq j$ means $F_i\subset F_j$, so indices higher
in this partial order correspond to deeper strata.  Besides
$i\succ\mathbf 0$,
the only strict chains consist of a two-stratum above one of its incident
four-strata or an isolated point above one of its two incident
four-strata, as displayed below:
\[
\text{two-stratum}\succ\text{four-stratum}\succ\mathbf 0,\qquad
\text{point}\succ\text{four-stratum}\succ\mathbf 0.
\]

We choose Joyce's dimension-indexed constants
$\zeta_0,\zeta_1,\zeta_2,\zeta_4,\zeta_8>0$ and define
$\xi_i=\zeta_{\dim V_i}$.  Only dimensions $0,2,4,8$ occur.  We set
\begin{equation}\label{eq:Joyce-neighborhoods}
 O_i=\{x:d(x,F_i)\leq\xi_i\},\qquad
 P_i=\{x:d(x,F_i)<2\xi_i\},\qquad
 Q_i=P_i\setminus\bigcup_{j\succ i}O_j.
\end{equation}
We choose the constants in increasing order of $\dim V_i$, beginning with
the zero-dimensional strata, and make each new constant sufficiently small
that two incomparable neighborhoods $P_i,P_j$ can meet only inside a
neighborhood $O_k$ of a component satisfying $k\succ i,j$.  Such a choice
exists because every intersection of two affine fixed flats is a finite
union of components belonging to deeper fixed sets.  Joyce's
\cite[Proposition 13.2.4]{JoyceBook} then proves
\[
 Q_i\cap Q_j\ne\varnothing\quad\Longrightarrow\quad
 i\succeq j\ \hbox{or}\ j\succeq i.
\]
The same proposition supplies, for each $i$, a $B_i$-equivariant flat
tubular identification of $Q_i/A_i$ with an open subset of
$F_i\times(W_i/A_i)$, and it proves that the $B_i$-action on that open set
is free.

For each singular index $i$, we denote by $\pi_i:Y_i\to W_i/A_i$ the model
appropriate to its stratum: an Eguchi--Hanson resolution in codimension
four, $Y_A$ in codimension six, and $\EH_\ell\times\EH_L$ at an isolated
point.  For $i=\mathbf 0$, we take the identity model.  We define
\begin{equation}\label{eq:Joyce-patches}
 M_i=(\id_{F_i}\times\pi_i)^{-1}(Q_i/A_i)
       \subset F_i\times Y_i.
\end{equation}
For $i\succeq j$, we denote by $\psi'_{i,j}$ the local-product embedding
constructed in the proof of Proposition~\ref{prop:R-data}, and we set
$\psi'_{i,i}=\id$.  We now define the relation in
\eqref{eq:Joyce-atlas} explicitly.  If $x\in M_j$ and $y\in M_k$, declare
$x\sim y$ precisely when one of the following two clauses holds:
\begin{enumerate}
\item There exists $\gamma\in K$ such that $i=\gamma\cdot k$ satisfies
$i\succeq j$, such that
$x\in M_j\cap(P_i\times Y_j)$, and
\begin{equation}\label{eq:atlas-relation-a}
 (\id\times\psi'_{i,j})(x)
   =(\gamma\times\chi_{\gamma,k})(y);
\end{equation}
\item There exists $\gamma\in K$ such that $i=\gamma\cdot j$ satisfies
$i\succeq k$, such that
$y\in M_k\cap(P_i\times Y_k)$, and
\begin{equation}\label{eq:atlas-relation-b}
 (\gamma\times\chi_{\gamma,j})(x)
   =(\id\times\psi'_{i,k})(y).
\end{equation}
\end{enumerate}
In both clauses, $\gamma\times\chi_{\gamma,k}:F_k\times Y_k\to
F_i\times Y_i$ acts by the affine map $\gamma:F_k\to F_i$ on the first
factor and by the resolution transport on the second.  These clauses
reproduce \cite[Definition 13.2.5]{JoyceBook}.  The cocycle
\eqref{eq:chi-cocycle} and equality of the incidence transition maps on
nested local-product charts imply reflexivity, symmetry, and transitivity
and identify the two composites on every triple overlap.  Hence
\begin{equation}\label{eq:Joyce-atlas}
 \widehat Z=\left(\coprod_{i\in\mathcal I}M_i\right)\big/\!\sim
\end{equation}
is covered by the smooth patches $M_i/B_i$.  The action of $B_i$ on $M_i$
is free because it covers the free action on $Q_i/A_i$.  Thus
\eqref{eq:Joyce-atlas} is a smooth
eight-manifold.  The maps $\id\times\pi_i$ glue to a blowdown
\begin{equation}\label{eq:Joyce-blowdown}
 \pi:\widehat Z\longrightarrow\T^8/K.
\end{equation}
Properness can be checked locally on the target.  We choose a finite
refinement of the Joyce cover by relatively compact open sets $U_\alpha$.
For each $\overline U_\alpha$, its inverse image under $\pi$ is a closed
subset of the union of finitely many inverse images under the proper model
maps $\id\times\pi_i$, and is therefore compact.  Thus $\pi$ is proper.
Since $\pi$ is surjective and $\T^8/K$ is compact, $\widehat Z$ is compact.
The atlas, equivalence relation, and blowdown just constructed are the
objects of \cite[Definition 13.2.5]{JoyceBook} in the present notation.

The scaled data entering the analytic theorem are as follows.  For
$t\in(0,1]$, we retain the same $Y_i$ and define
\[
 \pi_i^t=t\pi_i,\qquad
 \Phi_i^t=t^4D_{t;i}^*\Phi_i,\qquad
 g_i^t=t^2D_{t;i}^*g_i,
\]
where
\[
 D_{t;i}:V_i\times Y_i\longrightarrow V_i\times Y_i,
 \qquad D_{t;i}(v,y)=(t^{-1}v,y).
\]
Here $t\pi_i$ means postcomposition of $\pi_i$ with radial dilation by
$t$ on the cone $W_i/A_i$.
Thus the flat metric along $V_i$ is unchanged and the resolving normal
metric is scaled by $t^2$.  We define
\[
 M_i^t=(\id_{F_i}\times\pi_i^t)^{-1}(Q_i/A_i).
\]
For every strict incidence $i\succ j$, we define
\[
 V'_{ij}=V_j\cap W_i,\qquad V_j=V_i\oplus V'_{ij}.
\]
Condition (ii) gives an equivariant QALE identification
$\psi'_{i,j}:\mathcal U'_{ij}\to Y_i$ on an open overlap
$\mathcal U'_{ij}\subset V'_{ij}\times Y_j$.  We define
\[
 \mathcal U_{ij}=V_i\times\mathcal U'_{ij}
       \subset V_i\times V'_{ij}\times Y_j=V_j\times Y_j,
 \qquad
 \Psi_{ij}=\id_{V_i}\times\psi'_{i,j}.
\]
The displayed product decomposition identifies $\mathcal U_{ij}$ with an
eight-dimensional open subset of $V_j\times Y_j$.
Proposition~13.1.9 of \cite{JoyceBook} supplies an equivariant
$\sigma_{ij}\in\Omega^3(\mathcal U_{ij})$ satisfying
\begin{equation}\label{eq:unscaled-overlap-primitive}
 d\sigma_{ij}=\Psi_{ij}^*\Phi_i-\Phi_j.
\end{equation}
We define
\[
 D'_{t;ij}:V'_{ij}\times Y_j\longrightarrow V'_{ij}\times Y_j,
 \qquad D'_{t;ij}(v',y)=(t^{-1}v',y),
\]
\[
 \mathcal U_{ij}^{\prime,t}
   =(D'_{t;ij})^{-1}(\mathcal U'_{ij}),\qquad
\mathcal U_{ij}^t=D_{t;j}^{-1}(\mathcal U_{ij})
   =V_i\times\mathcal U_{ij}^{\prime,t},
\]
and we define
\[
 (\psi'_{i,j})^t
   =\psi'_{i,j}\circ D'_{t;ij}:
      \mathcal U_{ij}^{\prime,t}\longrightarrow Y_i,\qquad
 \Psi_{ij}^t=\id_{V_i}\times(\psi'_{i,j})^t
 =D_{t;i}^{-1}\circ\Psi_{ij}\circ D_{t;j},
 \qquad
 \sigma_{ij}^t=t^4D_{t;j}^*\sigma_{ij}.
\]
We use these scaled local-product maps and the unchanged equivariant
transports $\chi^t_{\gamma,i}=\chi_{\gamma,i}$ in the two clauses
\eqref{eq:atlas-relation-a}--\eqref{eq:atlas-relation-b}, with $M_i$
replaced by $M_i^t$.  The quotient defined by those clauses is denoted by
$\widehat Z^t$.  The local maps $\id\times\pi_i^t$ agree under the
equivalence relation and therefore define
\begin{equation}\label{eq:scaled-blowdown}
 \pi^t:\widehat Z^t\longrightarrow\T^8/K.
\end{equation}
The local models and all transition maps depend smoothly on $t$.  The
parameter-dependent patch charts and the graphs of their transition maps
therefore give the total quotient over $(0,1]$ the structure of a smooth
manifold, and projection to $t$ is a submersion with fibre $\widehat Z^t$.
Over a compact subinterval $J\subset(0,1]$, the properness argument used
for \eqref{eq:Joyce-blowdown}, applied to the finite parameterized patch
cover, shows that the inverse image of $J$ is compact.  Ehresmann's theorem
then makes the restriction over $J$ a smooth fibre bundle.  These
restrictions are compatible as $J$ varies, so the total quotient is a
smooth fibre bundle over $(0,1]$.  Since this interval is contractible,
the bundle admits a smooth trivialization.  We fix one such trivialization
$\vartheta_t:\widehat Z\to\widehat Z^t$ with $\vartheta_1=\id$ and pull
all subsequent forms and maps back by $\vartheta_t$.  We retain the same
symbols for the transported objects and also write $M_i^t$ for the
inverse-image patches in the fixed manifold $\widehat Z$.  This
gives the smooth connected family of scaled resolutions in
\cite[Definition 13.4.1]{JoyceBook}, written here with an explicit fixed
fibre.

The conjugation formula for $\Psi_{ij}^t$ is forced by the scaled
blowdowns.  Under the fixed identification
$W_i/A_i\cong V'_{ij}\times(W_j/A_j)$, it gives
\[
 \pi_i^t\bigl((\psi'_{i,j})^t(v',y)\bigr)
   =\bigl(v',\pi_j^t(y)\bigr).
\]
Applying $t^4D_{t;j}^*$ to
\eqref{eq:unscaled-overlap-primitive}, and using the conjugation formula,
gives
\begin{equation}\label{eq:overlap-primitive}
 d\sigma_{ij}^t=
 (\Psi_{ij}^t)^*\Phi_i^t-\Phi_j^t,
\end{equation}
and the scaled forms retain the full derivative estimates in
Proposition~13.1.9 of \cite{JoyceBook}.  We choose the unscaled primitives
simultaneously, as in \cite[Section 13.3.1]{JoyceBook}, so that
\[
(\Psi_{jk})^*\sigma_{ij}
 +\sigma_{jk}-\sigma_{ik}=0
\qquad(i\succ j\succ k),
\]
on every triple overlap.  Scaling this identity gives
\[
(\Psi_{jk}^t)^*\sigma_{ij}^t
 +\sigma_{jk}^t-\sigma_{ik}^t=0
\qquad(i\succ j\succ k),
\]
which is the additive cocycle needed on the scaled triple overlaps.

We now specify the correction terms.  For every four-dimensional
component $F_j$, Section 13.3.2 of \cite{JoyceBook} defines a homogeneous
three-form $\tau_j$ on $\R^8\setminus V_j$ as the leading homogeneous
term of $\sigma_{j\mathbf 0}$.  If $i\succeq j$, the $R$-data overlap map and the
flat tubular quotient push $\tau_j$ down to an invariant three-form on
$Q_i\setminus F_j$.  The pushed-down forms obtained from two comparable
indices coincide on their common domain.  We choose the special partition
of unity in \cite[Definition 13.2.6]{JoyceBook} and denote its functions
by $\eta_i'$.  These functions satisfy, simultaneously,
\[
 \operatorname{supp}\eta_i'\subset Q_i,\qquad
 \eta_{\gamma i}'\circ\gamma=\eta_i',\qquad
 \sum_{i\in\mathcal I}\eta_i'=1,
\]
and the following normal-direction condition.  For every
$i\ne\mathbf 0$ and
every $j$, there is a smooth function $\bar\eta_{ij}$ on $F_i$ such that,
on a neighborhood of $F_i$,
\[
 \eta_j'=\operatorname{pr}_{F_i}^*\bar\eta_{ij}.
\]
Thus every partition function, not only $\eta_i'$, is constant in the
$W_i$-directions there.  For a fixed index $j$ of dimension four, we
define, initially on the complement of $F_j$,
\begin{equation}\label{eq:Xi-smooth-extension}
 \Xi_j=\sum_{i\succeq j}d\eta_i'\wedge\tau_j
 =d\!\left(\sum_{i\succeq j}\eta_i'\right)\wedge\tau_j.
\end{equation}
Near $F_j$, every partition function that can be nonzero has an index
$i\succeq j$.  Consequently the parenthesized sum in
\eqref{eq:Xi-smooth-extension} is identically one on a neighborhood of
$F_j$, and its differential is identically zero there.  Therefore
$\Xi_j$ is zero on a punctured neighborhood of $F_j$ and has a unique
smooth extension across $F_j$ obtained by setting it equal to zero.  The
pushed-down forms agree on overlaps, so
\[
 \Xi=\sum_{\dim V_j=4}\Xi_j
\]
is a globally defined smooth $K$-invariant four-form on $\T^8$.  We use
the same symbol for the descended orbifold four-form on $\T^8/K$.
The flat metric and its Hodge star are $K$-invariant, so the Hodge
decomposition restricts to the $K$-invariant de Rham complex of
$\T^8$, equivalently to the orbifold de Rham complex of $\T^8/K$.
Applied to $d\Xi$, this decomposition first gives an invariant
anti-self-dual four-form $\Phi^-_0$ satisfying
$d\Phi^-_0=d\Xi$.  Any two such anti-self-dual solutions differ by a
harmonic anti-self-dual four-form.  The difference $\Phi^-_0-\Xi$ is
closed.  There is therefore a unique harmonic anti-self-dual four-form
$h^-$ for which
\[
 [\Phi^-_0+h^- -\Xi]\in H^4_+(\T^8/K;\R).
\]
Indeed, $h^-$ is the negative of the anti-self-dual harmonic component
of $[\Phi^-_0-\Xi]$.  We set $\Phi^-=\Phi^-_0+h^-$.  This proves both
existence and uniqueness of the anti-self-dual solution for which
$[\Phi^--\Xi]$ is self-dual.  We denote the harmonic self-dual
representative of this class by $\Psi$.  The closed form
$\Phi^--\Xi-\Psi$ represents zero in cohomology and is therefore exact.
We choose an invariant three-form $\Upsilon$ satisfying
\begin{equation}\label{eq:Joyce-global-correction}
 d\Upsilon=\Phi^--\Xi-\Psi.
\end{equation}
These are exactly the forms constructed in
\cite[Section 13.3.3]{JoyceBook}.

For an orbit $[i]\in\mathcal I/K$, we define
$\eta_{[i]}'=\sum_{j\in[i]}\eta_j'$.  This function is $K$-invariant and
therefore descends to $\T^8/K$.  We then define
\[
 \eta_{[i]}^t=(\pi^t)^*\eta_{[i]}'.
\]
The normal-direction condition implies that this pullback is smooth
across every exceptional set.  The functions $\eta_{[i]}^t$ form a
partition of unity on the fixed manifold $\widehat Z$.  We fix an index
$i$.  In a summand indexed by an orbit $[j]$, the relation $j\succ i$
(respectively, $i\succ j$) means that one has chosen a representative
$j\in[j]$ comparable with $i$.  The $K$-action and the
cocycle \eqref{eq:chi-cocycle} identify the local forms obtained from
any two such representatives.  Consequently each orbit sum below is
independent of this choice and descends to $M_i^t/B_i$.  On that quotient
patch, we define the pregluing form by
\begin{equation}\label{eq:Joyce-pregluing}
 \xi^t=\Phi_i^t
 +\sum_{\substack{[j]\in\mathcal I/K\\ j\succ i}}
       d(\eta_{[j]}^t\sigma_{ji}^t)
 -\sum_{\substack{[j]\in\mathcal I/K\\ i\succ j}}
       d(\eta_{[j]}^t\sigma_{ij}^t)
 +t^4d\bigl(f^t(\pi^t)^*\Upsilon\bigr)+t^4\Psi^t.
\end{equation}
The first two sums correct the exact differences
\eqref{eq:overlap-primitive}.  In the first sum every product
$\eta_{[j]}^t\sigma_{ji}^t$, and in the second every product
$\eta_{[j]}^t\sigma_{ij}^t$, is extended by zero outside the corresponding
orbit patch.  The support and normal-constancy conditions make both kinds
of extension smooth; equivariance gives the representative-independence
asserted above.

It remains to define $f^t$ and $\Psi^t$.  We denote the singular set by
$S\subset\T^8/K$ and the image of $F_i$ by $\overline F_i$.  We also define
\[
 \rho_t(x)=\bigl[t+d(\pi^t(x),S)\bigr]^{-1}.
\]
This is the weight function of
\cite[Definition 13.4.2]{JoyceBook}; the weighted estimates in
\cite[Propositions 13.4.3--13.5.7]{JoyceBook} are all taken with respect
to $\rho_t$.  Using the radii in \eqref{eq:Joyce-neighborhoods},
Definition 13.4.2 gives the open subsets
\begin{align*}
 D^t&=\{x:d(\pi^t(x),\overline F_i)>t^{1/3}\xi_i
             \text{ for every }i\ne\mathbf 0\},\\
 E^t&=\{x:d(\pi^t(x),\overline F_i)>2t^{1/3}\xi_i
             \text{ for every }i\ne\mathbf 0\}.
\end{align*}
The inclusion $E^t\subset D^t$ holds, and both sets lie over the regular
locus.  Definition 13.5.1 supplies a smooth function $f^t$ that is
identically zero on a neighborhood of $\widehat Z\setminus D^t$,
identically one on $E^t$, and satisfies
\[
 |\nabla^q f^t|_{g_i^t}=O(t^{-q/3})
 \quad\text{on }(M_i^t/B_i)\cap(D^t\setminus E^t),
 \qquad q=0,1,2,\qquad i\in\mathcal I.
\]
The form
$f^t(\pi^t)^*\Upsilon$ in \eqref{eq:Joyce-pregluing} is evaluated on the
regular support of $f^t$ and extended by zero; it is smooth because the
cutoff vanishes on a neighborhood of the complement of $D^t$.  The form
$\Psi^t$ is the closed extension furnished by
\cite[Proposition 13.4.4]{JoyceBook}; it equals $(\pi^t)^*\Psi$ on $D^t$ and
$[\Psi^t]=(\pi^t)^*[\Psi]$.  Equation
\eqref{eq:Joyce-global-correction} cancels the leading homogeneous
$t^4$ contribution from every four-dimensional stratum.  On an overlap,
equation \eqref{eq:overlap-primitive} and the additive cocycle for the
$\sigma_{ij}^t$ show that the two local expressions in
\eqref{eq:Joyce-pregluing} agree.  They therefore define a global
four-form $\xi^t$.  Each local expression is closed because
$d\Phi_i^t=0$, $d^2=0$, and $d\Psi^t=0$.  Hence $d\xi^t=0$.  This is
\cite[Definition 13.5.1 and Proposition 13.5.2]{JoyceBook}.

Let $\mathcal A\widehat Z\subset\Lambda^4T^*\widehat Z$ be the bundle of
admissible four-forms.  Definition 10.5.8 of \cite{JoyceBook} constructs
an open tubular neighborhood
$\mathcal T\widehat Z\subset\Lambda^4T^*\widehat Z$ and a specific smooth
fibre map $\Theta:\mathcal T\widehat Z\to\mathcal A\widehat Z$.  Pointwise,
each element of $\mathcal T\widehat Z$ has a unique expression
$\Omega+\chi$, where $\Omega$ is admissible,
$\chi\in\Lambda^4_{27}$ for the metric induced by $\Omega$, and $|\chi|$
is smaller than the fixed tubular radius; Joyce defines
$\Theta(\Omega+\chi)=\Omega$.  To avoid confusion with the holomorphic
volume form used earlier, we denote this particular map by
$\Theta_{\rm adm}$ below.  Proposition 10.5.9 gives its nonlinear
expansion and estimates.
For all sufficiently small $t$, the form $\xi^t$ takes values in the
domain of $\Theta_{\rm adm}$.  We define
\begin{equation}\label{eq:pre-torsion}
 \Phi_{\rm pre}^t=\Theta_{\rm adm}(\xi^t),\qquad
 \phi^t=\xi^t-\Phi_{\rm pre}^t,\qquad
 g^t=g_{\Phi_{\rm pre}^t}.
\end{equation}
Then
\begin{equation}\label{eq:torsion-primitive}
 d\Phi_{\rm pre}^t+d\phi^t=0.
\end{equation}
The cancellation in \eqref{eq:Joyce-pregluing} removes the leading
codimension-four torsion term.  The resulting estimates in
\cite[Theorem 13.5.8]{JoyceBook} give constants independent of $t$ such
that
\begin{equation}\label{eq:Joyce-estimates}
 \|\phi^t\|_{L^2}\le Ct^{13/3},\qquad
 \|d\phi^t\|_{L^{10}}\le Ct^{7/5},\qquad
 \delta(g^t)\ge ct,\qquad
 \|R(g^t)\|_{C^0}\le Ct^{-2}.
\end{equation}

\begin{theorem}\label{thm:Joyce-existence}
There exists $\varepsilon>0$ such that the compact smooth eight-manifold
$\widehat Z$ carries a torsion-free $\SpinSeven$ structure
$\widehat\Phi_t$ for every $t\in(0,\varepsilon)$.
\end{theorem}
\begin{proof}
Equations \eqref{eq:torsion-primitive}--\eqref{eq:Joyce-estimates} are
exactly the hypotheses of \cite[Theorem 13.6.1]{JoyceBook}.  In the notation
of its proof, there is a small anti-self-dual four-form $\eta^t$ satisfying
\[
 d\eta^t=d\phi^t+dF_{\Phi_{\rm pre}^t}(\eta^t),
 \qquad \|\eta^t\|_{C^0}=O(t^{1/3}),
\]
where the quadratic normal correction is characterized by
\[
 \Theta_{\rm adm}(\Phi_{\rm pre}^t+\eta)
 =\Phi_{\rm pre}^t+\eta-F_{\Phi_{\rm pre}^t}(\eta).
\]
Consequently,
\[
 \widehat\Phi_t=
 \Theta_{\rm adm}(\Phi_{\rm pre}^t+\eta^t)
 =\Phi_{\rm pre}^t+\eta^t-F_{\Phi_{\rm pre}^t}(\eta^t)
\]
is an admissible four-form.  Equations
\eqref{eq:torsion-primitive} and the nonlinear equation give the explicit
calculation
\begin{align*}
d\widehat\Phi_t
 &=d\Phi_{\rm pre}^t+d\eta^t
   -dF_{\Phi_{\rm pre}^t}(\eta^t)\\
 &=-d\phi^t+
   \bigl(d\phi^t+dF_{\Phi_{\rm pre}^t}(\eta^t)\bigr)
   -dF_{\Phi_{\rm pre}^t}(\eta^t)=0.
\end{align*}
An admissible four-form defines a $\SpinSeven$ structure, and such a
structure is torsion-free if and only if its defining four-form is closed.
Thus $\widehat\Phi_t$ is torsion-free.  This is the existence
conclusion of \cite[Theorem 13.6.2]{JoyceBook}.
\end{proof}

\section{The equivariant first resolution and the global blowdown}
The Joyce resolution admits a topological description by two successive
fibrewise resolutions.  This description identifies its cohomology and
fundamental group in Sections~5 and~6.

\subsection{The Mart\'in--Merch\'an resolution}
The singular locus of $X$ is the disjoint union of ten associative
three-manifolds $N_1,\ldots,N_{10}$ of exact normal isotropy $\Z_2$
\cite[Section 2]{MM}.  Using the component frames, radii, and lifted
transition maps already fixed for the old $G$-strata in the Joyce atlas,
we denote the corresponding first-stage resolution by
\begin{equation}\label{eq:M-resolution}r:M\longrightarrow X.\end{equation}
The relative equivariant
tubular-neighborhood argument at the end of Lemma~\ref{lem:tau-lift}
identifies it, over $X$ and relative to the exterior, with the resolution
used in \cite{MM}.  It replaces each normal
$\C^2/\{\pm1\}$ by
\begin{equation}\label{eq:Y-model}
Y=\{(w,\ell)\in\C^2\times\CP^1:w\in\ell\}/((w,\ell)\sim(-w,\ell)).
\end{equation}
The fibrewise squaring map identifies this smooth complex surface with
$\operatorname{Tot}\mathcal O_{\CP^1}(-2)$, the minimal resolution of
$\C^2/\{\pm1\}$.
The manifold $M$ is simply connected \cite[Section 2.2]{MM}.  We write
$E_j=r^{-1}(N_j)$ and denote its exceptional Thom class by
$x_j\in H^2(M)$.  The normal complex
structure is defined by
\begin{equation}\label{eq:normal-I}
g(I_jv,w)=\varphi(\chi_j,v,w),\qquad
\chi_j=\begin{cases}\partial_{y_3},&j=1,2,\\
\partial_{x_3},&j\ge3.\end{cases}
\end{equation}

\begin{lemma}\label{lem:tau-lift}
The involution $\tau$ lifts to an involution $\widetilde\tau:M\to M$ with
$r\circ\widetilde\tau=\tau\circ r$, and
\begin{equation}\label{eq:tau-x}
\widetilde\tau^*x_1=-x_2,\qquad \widetilde\tau^*x_2=-x_1,\qquad
\widetilde\tau^*x_j=x_j\quad(3\le j\le10).
\end{equation}
\end{lemma}
\begin{proof}
At $t=0$, we express the action in the same mapping-torus chart by defining
\begin{equation}\label{eq:tau0}
\tau_0=F^{-1}\tau,\qquad
\tau_0(z_1,z_2,z_3)=(-\ii z_1,-\ii z_2,z_3-\tfrac14).
\end{equation}
At $t=0$ the two relevant chosen covers are
\begin{align*}
N_1^1&=\{(x_1-\ii x_1,x_2-\ii x_2,\ii y_3):x_1,x_2,y_3\in\R/\Z\},\\
N_2^1&=\{(x_1+\ii x_1,x_2+\ii x_2,-\tfrac14+\ii y_3):
             x_1,x_2,y_3\in\R/\Z\}.
\end{align*}
Consequently $\tau_0$ sends the displayed point of $N_1^1$ to
\[
(-x_1-\ii x_1,-x_2-\ii x_2,-\tfrac14+\ii y_3)\in N_2^1.
\]
The same substitution on $N_2^1$ gives
\[
\tau_0(N_1^1)=N_2^1,\qquad
\tau_0(N_2^1)=N_1^2=\iota_1N_1^1.
\]
At $t=\frac12$, the covers of $N_3,\ldots,N_6$ have fixed imaginary
half-periods in the three elliptic coordinates, and $\tau$ preserves each
chosen component.  The covers of $N_7,\ldots,N_{10}$ have fixed real
half-periods in $z_1,z_2$ and a fixed imaginary quarter-period in $z_3$;
$\iota_2\tau$ preserves each chosen component.  Thus the maps on the
component covers are
\begin{equation}\label{eq:component-lifts}
\begin{array}{c|c|c}
\text{components}&\text{map on the component cover}&\text{deck involution}\\
\hline
N_1^1\longrightarrow N_2^1&\tau_0&\iota_2\\
N_2^1\longrightarrow N_1^1&\iota_1\tau_0&\iota_2\\
N_j^1\longrightarrow N_j^1,\ 3\le j\le6&\tau&\iota_1\\
N_j^1\longrightarrow N_j^1,\ 7\le j\le10&\iota_2\tau&\iota_1.
\end{array}
\end{equation}
Here $\tau_0$ commutes with $\iota_1,\iota_2$ and
$\tau_0^2=\iota_1$.  Thus the first two displayed maps are inverse to one
another and are $\iota_2$-equivariant.  The maps in the last two rows square
to the identity and commute with $\iota_1$.

Every representative $L$ in \eqref{eq:component-lifts} satisfies
$L^*\varphi=-\varphi$.  If $L_*\chi_j=\chi_k$, then for normal vectors
$v,w$,
\[
g(I_kLv,Lw)=\varphi(\chi_k,Lv,Lw)
           =-\varphi(\chi_j,v,w)=g(-LI_jv,Lw).
\]
Thus $I_kL=-LI_j$.  In particular, $D\tau_0\chi_1=\chi_2$ gives
\begin{equation}\label{eq:antilinear}D\tau_0I_1=-I_2D\tau_0.\end{equation}
The same identity holds for $\iota_1\tau_0$.  In the last two rows of
\eqref{eq:component-lifts}, the derivative sends $\chi_j$ to $-\chi_j$;
the same displayed calculation now has two minus signs, which cancel, and
\begin{equation}\label{eq:linear}D\tau I_j=I_jD\tau\quad(j\ge3).\end{equation}
Here $D\tau$ denotes the derivative of the representative in the relevant
row; the translations $\iota_2$ do not alter the calculation.

We now construct the global lift, including the normal orbibundle
quotients.  In the notation of Mart\'in--Merch\'an, the resolved local piece
over $N_j$ is
\begin{equation}\label{eq:MM-local-piece}
 \mathcal P_j=(N_j^1\times Y)/\langle d_j\rangle,
 \qquad
 d_j=\begin{cases}\iota_2,&j=1,2,\\
                   \iota_1,&3\leq j\leq10.
       \end{cases}
\end{equation}
In the unitary normal frame used in \eqref{eq:normal-I}, the deck map acts
on the two complex normal coordinates by $\diag(1,-1)$; the central
$-I$ has already been divided out in the definition of $Y$.  Thus
\[
 d_j(x,[(w,\lambda)])=
 (d_jx,[(D d_jw,Dd_j\lambda)]).
\]
A linear or antilinear isometry $L$ lifts canonically to $Y$ by
\begin{equation}\label{eq:canonical-Y-lift}
 [(w,\lambda)]\longmapsto[(Lw,L\lambda)].
\end{equation}

We define $L_1=\tau_0:N_1^1\to N_2^1$ and
$L_2=\iota_1\tau_0:N_2^1\to N_1^1$.  We also define
$L_j=\tau$ for $3\leq j\leq6$ and $L_j=\iota_2\tau$ for
$7\leq j\leq10$.  The commutation relations in
\eqref{eq:component-lifts} hold on both the base and the normal frame, so
\begin{equation}\label{eq:deck-equivariance}
 L_jd_j=d_kL_j
\end{equation}
when $L_j$ maps the $j$th component to the $k$th.  Hence
\[
 \widehat L_j[x,y]=[L_jx,\widehat{D L_j}y]
\]
is a well-defined map between the quotients \eqref{eq:MM-local-piece}.
The compositions are exact:
\[
 (\iota_1\tau_0)\tau_0=\iota_1\tau_0^2=1,
 \qquad \tau_0(\iota_1\tau_0)=1,
 \qquad \tau^2=(\iota_2\tau)^2=1.
\]
Thus the lifted local maps satisfy the involution cocycle before gluing.

For the tubular gluing, we choose an ambient metric on a neighborhood of
the disjoint union of the lifted strata in the smooth mapping-torus cover
and average it over the finite group generated by the old deck maps and
the $L_j$.  The resulting orthogonal normal bundles and ambient
exponential maps are equivariant.  We choose the radii and radial cutoff
functions to be constant on each orbit of components.
On the common annulus, \eqref{eq:canonical-Y-lift} agrees with the original
linear action through the blowdown, so it glues to that action on the
complement.  The maps $\widehat L_j$ therefore define a global involution
$\widetilde\tau$ with $r\circ\widetilde\tau=\tau\circ r$.

Finally, this equivariant choice has the same underlying resolution as the
one used in \cite{MM}.  Indeed, the equivariant tubular-neighborhood
uniqueness theorem \cite[Chapter VI]{BredonGroup} gives an equivariant
isotopy between any two such tubular embeddings after shrinking the radii;
isotopy extension \cite[Chapter 8]{Hirsch} gives a
diffeomorphism relative to the complement of the resolution
neighborhoods.  Hence the exceptional labels and the cohomology-ring
calculation from \cite{MM} apply to the present $r$.  For $j=1,2$ the lift
is antiunitary and conjugates the $\mathcal O(-2)$ fibre coordinate, so it
reverses the Thom orientation.  For $j\geq3$ it is unitary and preserves
that orientation.  This proves \eqref{eq:tau-x}.
\end{proof}

The sign in \eqref{eq:epsilon} is equivalently determined by
$D\gamma I_C=\epsilon_{\gamma,C}I_{\gamma C}D\gamma$.  Thus its cocycle law
is automatic.  In particular $\epsilon_{\tau,N_1}=-1$, while it is $+1$ on
$N_j$, $j\ge3$; the $R$-data uses the same exceptional orientations.

\subsection{The second stage}
We define
\begin{equation}\label{eq:W-Z0}
W=S^1\times M,\qquad \sigma(s,p)=(-s,\widetilde\tau p),\qquad
Z_0=W/\langle\sigma\rangle.
\end{equation}
Because $\tau$ normalizes $G$, it induces $\bar\tau:X\to X$, and there is a
canonical identification
\[
\mathcal O=(S^1\times\T^7)/(G\rtimes\langle T\rangle)
\cong(S^1\times X)/\langle\bar T\rangle,
\qquad \bar T(s,x)=(-s,\bar\tau x).
\]
The identity
$(\id_{S^1}\times r)\circ\sigma
=\bar T\circ(\id_{S^1}\times r)$ makes
$\id_{S^1}\times r$ descend to a blowdown
\begin{equation}\label{eq:q0}
q_0:Z_0\longrightarrow\mathcal O,\qquad
q_0([s,p])=[s,r(p)].
\end{equation}
It is proper because its domain is compact.

\begin{lemma}\label{lem:two-stage}
The fixed locus of $\sigma$ is a disjoint union of smooth four-manifolds.
Resolving the resulting $A_1$ singularities gives a smooth manifold
$\widehat Z_{\rm st}$ and a proper degree-one map
\begin{equation}\label{eq:rho}\rho:\widehat Z_{\rm st}\longrightarrow Z_0.\end{equation}
This staged resolution is diffeomorphic, compatibly over the regular locus,
to the Joyce resolution $\widehat Z$.
\end{lemma}
\begin{proof}
Since $\sigma$ is a smooth involution, its fixed locus is a smooth embedded
submanifold.  Away from the old exceptional set, the exact-isotropy part of
Table~\ref{tab:coset-fixed} says that all its branches have dimension four.
Every zero-dimensional branch lies on an old stratum by
\eqref{eq:zero-old-stabilizer}.  There are two local intersection models.

At a codimension-six stratum, we choose coordinates in which
\begin{align}\label{eq:stage-SU3}
a(u,z_1,z_2,z_3)&=(u,z_1,-z_2,-z_3),\\
h(u,z_1,z_2,z_3)&=(u,-z_1,z_2,-z_3).\nonumber
\end{align}
We resolve the old $a$-quotient in $(z_2,z_3)$ by the scale-$L$
Eguchi--Hanson model $Y$.  The residual map is
induced by
\[
(z_1,z_2,z_3)\longmapsto(-z_1,z_2,-z_3).
\]
On $Y$, the lift of
$(z_2,z_3)\mapsto(z_2,-z_3)$ has two disjoint fixed complex lines, over
$[1:0]$ and $[0:1]$ in the exceptional $\CP^1$.  More explicitly, if
\[
L_2=\{[((w_2,0),[1:0])]:w_2\in\C\},\qquad
L_3=\{[((0,w_3),[0:1])]:w_3\in\C\},
\]
then the residual fixed set is
\[
\R^2_u\times\{0\}_{z_1}\times(L_2\sqcup L_3).
\]
It consists of two disjoint real four-dimensional branches.  We resolve
both branches.  This operation inserts the remaining two midpoint rays
with scale $\ell$ and gives the staged fan
\eqref{eq:toric-cones} and the support data \eqref{eq:support-values}.

At an isolated product point, we choose the representation
$a=(I_U,-I_{U'})$ and $h=(-I_U,-I_{U'})$.
After resolving the $a$-quotient, the model is
\[
 U\times\EH_L(U'),
\]
and the residual involution is
$(-I_U,\id_{\EH_L(U')})$.  Its fixed set is
$\{0\}\times\EH_L(U')$, again of real dimension four.  The second
resolution is
$\EH_\ell(U)\times\EH_L(U')$ as in \eqref{eq:EH-product}.

These models are exhaustive for the following group-theoretic reason.  If
$p\in\Fix(\sigma)$ lies over the old exceptional set, then
$q_0([p])\in\mathcal O$ has stabilizer containing a nonidentity old element
$a\in G$ and a new element $h\in GT$.  Proposition~\ref{prop:isotropy}
says that $\langle a,h\rangle$ is either a codimension-six $\Z_2^2$ with
fixed-dimension pattern $(4,4,4)$, or an isolated $\Z_2^2$ with pattern
$(4,4,0)$; no larger stabilizer occurs.  The first case is exactly the
$L_2\sqcup L_3$ calculation and the second is exactly
$\{0\}\times\EH_L(U')$.  Points away from the old exceptional set have
exact new $\Z_2$ isotropy and give the regular four-dimensional branches
already identified from Table~\ref{tab:coset-fixed}.  Thus every component
$B$ of $\Fix(\sigma)$ has codimension four.

We now define the fibrewise resolution globally.  For each component $B$
of $\Fix(\sigma)$, we identify $B$ with its image in $Z_0$ and write
$\nu B=TW|_B/TB$ for its normal bundle.  The derivative of $\sigma$ is the
identity on $TB$ and minus the identity on $\nu B$.

On the dense regular locus, $\widetilde\tau$ agrees with $\tau$, whose
linear part has determinant $-1$; hence $\widetilde\tau$ reverses the
orientation of $M$.  Reflection reverses the orientation of $S^1$, so
$\sigma$ preserves the product orientation of $W$.  Each tangent four-plane of $B$ has the
Cayley orientation in its local product model, and the transition maps
preserve the corresponding Cayley form.  These local orientations give a
global orientation of $TB$.  We orient $\nu B$ so that the direct-sum
orientation of $TB\oplus\nu B$ is the ambient orientation.  In
each $R$-data chart, the chosen normal frame identifies the normal fibre
with the standard oriented $\R^4=\C^2$.  Let
\[
 H_B=U(2)\rtimes\langle\text{complex conjugation}\rangle\subset SO(4).
\]
The normalizer matrices \eqref{eq:centralizer-image} and
\eqref{eq:two-normalizer-signs}, together with the cocycle
\eqref{eq:chi-cocycle}, show that every change between these normal frames
takes values in $H_B$ and that the changes satisfy the cocycle identity.
In a codimension-six chart, the two residual branches are $L_2,L_3$, and
\eqref{eq:incidence-signs} identifies the normal-frame reduction obtained
from the regular chart with the one obtained from the
$\C\times\EH_L$ chart.  At an isolated chart, the two factors in
\eqref{eq:EH-product} give the same identification.  Consequently the
normal frame cocycle defines a principal $H_B$-reduction $P_B$ of the
oriented linear frame bundle of $\nu B$.

The standard inner product on $\R^4$ is $H_B$-invariant, so $P_B$ defines
a metric on $\nu B$.  The standard K\"ahler form is self-dual for this
metric; $U(2)$ fixes it and complex conjugation negates it.  Hence its
unoriented span defines a smooth line bundle
\[
 \ell_B=P_B\times_{H_B}\R\omega_{\rm std}
 \subset\Lambda^+\nu^*B.
\]
This order of construction is essential: the $H_B$-reduction first
defines the compatible normal metric, and self-duality is asserted only
for that metric.  We extend the metric from $TB\oplus\nu B$ to a
$\sigma$-invariant metric near $B$ and then to $W$; averaging does not
alter its prescribed value along $B$ because $D\sigma$ acts as
$I\oplus(-I)$.  The equivariant exponential map now identifies a
neighborhood of $B$ in $Z_0$ with the orbibundle
\begin{equation}\label{eq:normal-orbibundle}
 \nu B/\{\pm1\}\longrightarrow B,
\end{equation}
whose fibre is $\R^4/\{\pm1\}$.

Every transition in $H_B$ has the canonical radial
holomorphic or antiholomorphic lift to $\EH_\ell$ constructed in
\eqref{eq:epsilon}; equation \eqref{eq:chi-cocycle} says that these lifted
transition functions satisfy the cocycle.  Using the principal
$H_B$-reduction $P_B$, we define
\begin{equation}\label{eq:associated-EH-bundle}
 \mathcal E_B=P_B\times_{H_B}\EH_\ell,
 \qquad
 \rho_B([f,y])=[f,\pi_{\EH}(y)]
   \in P_B\times_{H_B}(\R^4/\{\pm1\}).
\end{equation}
Here the notation means that $H_B$ acts through its canonical lifted
action on $\EH_\ell$; equivalently, \eqref{eq:associated-EH-bundle} is
defined directly by the lifted transition cocycle.  We truncate both
bundles at the common chosen radius.  On the boundary annulus the Eguchi--Hanson
blowdown is the original quotient identification, so $\mathcal E_B$ glues
to the unchanged exterior by the same fixed boundary identification.  Performing
this construction on the disjoint components defines the global map
$\rho$ in \eqref{eq:rho}.  It is an orientation-preserving diffeomorphism
over the regular set; a regular value has one positive preimage, so
$\deg\rho=1$.

It remains to compare the gluing with the $R$-data resolution.  The map
$r$ was defined using the restrictions of the Joyce component frames,
tubular maps, radii, and lifted cocycle to the old subgroup $G$.  At a
two-stratum, the first blowup is the toric star subdivision by $m_{23}$ in
$N_a$.  Passing to the residual quotient enlarges $N_a$ to $N$; the two
nonsaturated boundary faces are $\langle e_1,e_2\rangle$ and
$\langle e_1,e_3\rangle$, and their star subdivisions insert $m_{12}$ and
$m_{13}$.  Hence the resulting fan is exactly \eqref{eq:toric-cones},
including its diagonal $\langle e_1,m_{23}\rangle$ and the support values
corresponding to periods $L,\ell,\ell$.  Every transition derivative lies in the relevant
normalizer after the labels are transported from $A_0$ or $A_1$.
The exact calculation \eqref{eq:two-normalizers}--\eqref{eq:two-normalizer-signs}
shows that it fixes every ray label; in particular it preserves the fan and
support function.  Its toric lift is exactly the lift used in the $R$-data
atlas.  At a four-stratum the
analogous statement is uniqueness of the lift to the minimal $A_1$
resolution, and at a zero-stratum it is the product lift to
$\EH_\ell(U)\times\EH_L(U')$.  Thus, for this chosen representative of the
first-stage resolution, the staged atlas and the $R$-data atlas have the
same local models, tubular maps, and lifted transition cocycle.  In a
codimension-six
chart the composite ``$m_{23}$ subdivision, residual lattice quotient,
$m_{12},m_{13}$ subdivisions'' is the toric transition associated with
the fan \eqref{eq:toric-cones}; at a point it is the product blowdown
\eqref{eq:EH-product}.  On every boundary annulus both use the original
affine quotient map.  Hence the transition maps coincide and glue to a
global diffeomorphism over $\mathcal O$ whose restriction to the regular
locus is the identity.

This identification does not depend on the chosen presentation of the
first resolution.  After the tubular radii are decreased if necessary,
any two equivariant tubular embeddings of the same normal orbibundles are
equivariantly isotopic.  The isotopy lifts through the fixed local
resolution model and extends across the common boundary annulus.  If the
strata are treated in decreasing order of depth, the isotopy extension
theorem produces at each stage a diffeomorphism relative to the pieces
already treated and relative to the exterior.  This relative uniqueness
identifies the chosen map $r$ with the resolution map of any
Mart\'in--Merch\'an presentation while preserving the exceptional labels
and the cohomology classes used below.
\end{proof}

We henceforth identify $\widehat Z_{\rm st}$ with $\widehat Z$ by the
diffeomorphism constructed in Lemma~\ref{lem:two-stage}.

\section{A Massey product outside its full indeterminacy}
We define
\begin{equation}\label{eq:abc}
a=x_1+x_2,\qquad b=x_7+x_3,\qquad c=x_7-x_3\in H^2(M).
\end{equation}

We use the following set-valued convention.  If closed forms
$\alpha_1,\alpha_2,\alpha_3$ of degrees $p,q,r$ represent
$a_1,a_2,a_3$, and
\[
d\xi_{12}=\alpha_1\alpha_2,\qquad
d\xi_{23}=\alpha_2\alpha_3,
\]
then
\begin{equation}\label{eq:Massey-coset}
\langle a_1,a_2,a_3\rangle
=[\xi_{12}\alpha_3-(-1)^p\alpha_1\xi_{23}]
+I(a_1,a_3),
\end{equation}
where the full indeterminacy is
\begin{equation}\label{eq:general-indeterminacy}
I(a_1,a_3)=a_1H^{q+r-1}+H^{p+q-1}a_3
\end{equation}
in the cohomology of the ambient CDGA.
As the representatives and primitives vary, the resulting cohomology
classes form precisely the affine coset in \eqref{eq:Massey-coset}.
Consequently zero belongs to the Massey product if and only if one, and
hence every, defining-system representative belongs to the full
indeterminacy subspace \eqref{eq:general-indeterminacy}.

\begin{proposition}[Mart\'in--Merch\'an]\label{prop:MM-Massey}
The product $\langle a,b,c\rangle$ is defined and
\begin{equation}\label{eq:MM-nonzero}0\notin\langle a,b,c\rangle.\end{equation}
Equivalently, every defining-system representative $Y\in H^5(M)$ satisfies
\begin{equation}\label{eq:old-indeterminacy}
Y\notin I_M,\qquad I_M=aH^3(M)+H^3(M)c.
\end{equation}
\end{proposition}
\begin{proof}[Full-indeterminacy calculation]
The additive splitting and products from \cite[Section 3]{MM} are
\begin{align}
H^k(M)&=H^k(X)\oplus\bigoplus_{j=1}^{10}H^{k-2}(N_j)x_j,\label{eq:MM-cohom}\\
x_jx_k&=0\ (j\ne k),\qquad x_j^2=-2\PD[N_j].\label{eq:exceptional-products}
\end{align}
For completeness, these formulas follow directly from the local
replacement.  We denote by $U_0$ the complement of smaller disjoint tubular
neighborhoods of the $N_j$.  Both $X$ and $M$ are obtained from $U_0$ by
attaching bundles with the same boundary $N_j\widetilde\times\RP^3$; the
singular fibre is $C(\RP^3)$ and the resolved fibre is the truncated disc
bundle of $\mathcal O_{\CP^1}(-2)$.  Over $\R$, their fibre-pair
cohomologies are, explicitly,
\[
H^q(C(\RP^3),\RP^3;\R)=
 \begin{cases}\R,&q=4,\\0,&q\ne4,\end{cases}
\quad
H^q(D(\mathcal O(-2)),S(\mathcal O(-2));\R)=
 \begin{cases}\R,&q=2,4,\\0,&q\ne2,4.\end{cases}
\]
We spell out the bundle argument that turns this fibre calculation into a
calculation of the relative groups.  For each $j$, let $T_j$ be the relative
orientation class of the singular cone bundle, normalized to have fibre
integral one, and choose a closed, compactly supported Thom form $\theta_j$
of the exceptional divisor in the resolved bundle.  On a resolved fibre,
$[\theta_j]$ generates the degree-two relative cohomology and
$-\frac12[\theta_j]^2$ generates the degree-four relative cohomology,
because the exceptional curve has self-intersection $-2$.  The local
coefficient systems are trivial.  Indeed, the degree-four system is the
orientation system of the oriented normal cone, while in
\eqref{eq:MM-local-piece} the deck transformation acts holomorphically on
$\CP^1$ and therefore fixes both $[\theta_j]$ and
$[\theta_j]^2$.

Relative Leray--Hirsch, followed by excision, now gives
\begin{align*}
 H^k(X,U_0)&\cong\bigoplus_j H^{k-4}(N_j),\\
 H^k(M,U_0)&\cong\bigoplus_j\left(
 H^{k-2}(N_j)[\theta_j]\oplus H^{k-4}(N_j)[\theta_j]^2\right).
\end{align*}
The fibrewise blowdown has oriented degree one.  Therefore the relative
pullback of $T_j$ has fibre integral one, so its
$-\frac12[\theta_j]^2$-coordinate is one.  Consequently relative pullback
is injective, its image is a graph over the degree-four summand, and its
cokernel is identified, using the chosen Leray--Hirsch generators, with
$\bigoplus_jH^{k-2}(N_j)[\theta_j]$.  Choosing this graph as the common
summand gives the relative splitting used below.  We define
\begin{equation}\label{eq:cohomology-splitting-map}
 \Psi_k\bigl(\xi,(\lambda_j)_j\bigr)
 =r^*\xi+\sum_j
   [\operatorname{pr}_j^*\lambda_j\wedge\theta_j].
\end{equation}
Here $\xi$ and the $\lambda_j$ are cohomology classes.  To realize the
displayed map by differential forms, we choose a representative of $\xi$
that is constant in the normal directions near every $N_j$ before taking
its pullback, as in \cite[Section 3]{MM}, and we choose closed
representatives of the $\lambda_j$ on the $N_j$.

We now give the diagram chase that proves that $\Psi_k$ is an
isomorphism.  The vertical map on the common exterior $U_0$ is the
identity.  We fix $y\in H^k(M)$ and write $u=y|_{U_0}$.  The boundary of $u$
in the long exact sequence of $(M,U_0)$ is zero.  Naturality and
injectivity of relative pullback on the common
$\bigoplus_jH^{k+1-4}(N_j)$ summand imply that the boundary of $u$ in the
long exact sequence of $(X,U_0)$ is also zero.  Exactness therefore gives
a class $\xi\in H^k(X)$ satisfying $\xi|_{U_0}=u$.  The class
$y-r^*\xi$ restricts to zero on $U_0$ and hence is the image of a relative
class in $H^k(M,U_0)$.  Decompose that relative class into its common
degree-four part and its additional degree-two part.  The common part is
the pullback of a relative class for $(X,U_0)$ and can be absorbed into
$\xi$; the additional part maps to a sum of the supported classes in
\eqref{eq:cohomology-splitting-map}.  This proves that $\Psi_k$ is
surjective.

The new summands are independent of one another and of the pullback
summand.  Indeed, if
$i_j:E_j\hookrightarrow M$ and $p_j:E_j\to N_j$, then
\begin{equation}\label{eq:exceptional-coefficient-map}
 -\tfrac12(p_j)_*i_j^*
 \bigl(\operatorname{pr}_k^*\lambda_k\wedge\theta_k\bigr)
 =\delta_{jk}\lambda_k,
 \qquad
 (p_j)_*i_j^*r^*\xi=0.
\end{equation}
Indeed, $i_j^*\theta_j=e(\mathcal O(-2))$ integrates to $-2$ on the
$\CP^1$ fibre.  Thus the homomorphism
$C_j=-\frac12(p_j)_*i_j^*$ extracts the $j$th exceptional coefficient.

It remains to prove that $r^*:H^k(X)\to H^k(M)$ is injective.  Suppose
that $r^*\xi=0$.  Its restriction to $U_0$ is zero, so exactness represents
$\xi$ by a relative class $a\in H^k(X,U_0)$.  Naturality shows that the
relative pullback of $a$ maps to zero in $H^k(M)$.  Exactness for
$(M,U_0)$ therefore expresses that pullback as the boundary of a class in
$H^{k-1}(U_0)$.  The relative pullback is injective on the common relative
summand, so naturality implies that $a$ itself is the boundary of the same
class in the sequence for $(X,U_0)$.  Hence $\xi=0$.  The maps $C_j$ and
the injectivity of $r^*$ prove that $\Psi_k$ is injective.  Therefore every
$\Psi_k$ is an isomorphism, proving \eqref{eq:MM-cohom} with the stated
normalization.  Supports for
different resolution neighborhoods are
disjoint, so $x_jx_k=0$ for $j\ne k$.  On one fibre, the normal line of the
exceptional $\CP^1$ is $\mathcal O(-2)$; its Euler class evaluates to
$-2$.  The Thom self-intersection formula therefore gives
$x_j^2=-2\PD[N_j]$, including the displayed sign.  This also proves the
supported pushforward rule \eqref{eq:pushforward-product} used below.
Here $\PD[N_j]$ means the pullback of the Poincar\'e dual in $X$, regarded
as an element of the $H^4(X)$ summand in \eqref{eq:MM-cohom}.  We denote
the exceptional projection by
\begin{equation}\label{eq:piE}
\pi_E:H^5(M)\longrightarrow
\bigoplus_{j=1}^{10}H^3(N_j)x_j.
\end{equation}
The two products needed to define the Massey product are
\[
ab=(x_1+x_2)(x_7+x_3)=0
\]
because the four exceptional neighborhoods are pairwise disjoint, and
\[
bc=x_7^2-x_3^2=-2\PD[N_7]+2\PD[N_3]=0
\]
because $\PD[N_3]=\PD[N_7]$.  Thus
$\langle a,b,c\rangle$ is defined.  We define
$v_j=[\varphi|_{N_j}]x_j$.  We recall the geometric
construction of the primitive used by Mart\'in--Merch\'an.  On the
mapping-torus cover, $N_3^1$ and $N_7^1$ are joined by the embedded
oriented cobordism
\[
 \mathcal C=\{(t,x_1,x_2,x_3+\ii f(t)):
  \tfrac12\leq t\leq\tfrac32,\ x_1,x_2,x_3\in\R/\Z\},
\]
where $f$ is nondecreasing, $f=0$ for $t\leq9/8$, and $f=1/4$ for
$t\geq5/4$.  Two points in the displayed parameter domain have the same
image only if their $t$-coordinates differ by an integer.  Since the
interior $t$-coordinates lie in an interval of length one, equality in the
interior forces equality of all parameters.  At the endpoints, the only
additional identification is the defining mapping-torus identification.
It identifies the upper boundary with the specified lift of $N_7$ and
does not identify distinct interior points.  Hence the parametrization is
an embedding, and the boundary-orientation convention gives
$\partial\mathcal C=N_3^1-N_7^1$.  Its normal bundle is trivialized by
\[
 \partial_{y_1},\quad\partial_{y_2},\quad
 \partial_{y_3}-f'(t)\partial_t.
\]
We choose $\zeta:\R^3/\Z^3\to\R$ with integral one and support in
$B_{2\delta}(0)\setminus B_\delta(0)$.  We also choose $h$ to be zero near
$t=1/2,3/2$, equal to one on
$[1/2+2\epsilon,3/2-2\epsilon]$.  In the displayed normal coordinates, we define
\begin{equation}\label{eq:explicit-cobordism-form}
 \beta_{\mathcal C}=
 h(t)\zeta(y_1,y_2,y_3-f(t))\,dy_1\wedge dy_2
 \wedge\bigl(dy_3-f'(t)\,dt\bigr).
\end{equation}
This is the Thom form of $\nu\mathcal C$ away from the boundary collars.
We choose $h$ to be increasing on the lower boundary collar and decreasing
on the upper boundary collar.  We denote the four-form supported in the
lower collar by $\upsilon_3^1$ and the negative of the four-form supported
in the upper collar by $\upsilon_7^1$.  Then
\[
 d\beta_{\mathcal C}
 =h'(t)\zeta(y_1,y_2,y_3-f(t))\,
   dt\wedge dy_1\wedge dy_2\wedge dy_3
 =\upsilon_3^1-\upsilon_7^1.
\]
The normalization $\int\zeta=1$, together with
$\int h'(t)\,dt=1$ on the lower collar and
$\int h'(t)\,dt=-1$ on the upper collar, shows that both
$\upsilon_3^1$ and $\upsilon_7^1$ have fibre integral $+1$ with respect to
the oriented normal fibres of their respective boundary components.

We define $K_0=\langle\iota_1,\iota_2,\kappa\rangle$, which has order
eight, and set
\[
 \beta_0=\frac14\sum_{\gamma\in K_0}\gamma^*\beta_{\mathcal C},
 \qquad
 \upsilon_j^0=\frac14\sum_{\gamma\in K_0}\gamma^*\upsilon_j^1
 \quad(j=3,7).
\]
These forms are $K_0$-invariant and satisfy
$d\beta_0=\upsilon_3^0-\upsilon_7^0$.  The subgroup of $K_0$ preserving a
chosen lifted component has order four.  Thus, on a lifted normal fibre,
the four coincident stabilizer terms cancel the factor $1/4$, and
$\upsilon_j^0$ has fibre integral one.  The invariant forms descend
through the old quotients.  On a normal fibre of $N_j$, the remaining
quotient is the degree-two map
$\C^2\to\C^2/\{\pm1\}$.  Consequently the descended closed forms
$\upsilon_j$ on $X$ have the precise normalization
\begin{equation}\label{eq:half-Thom}
 2[\upsilon_j]=\operatorname{Th}[N_j]\in H_c^4(V_j),
 \qquad j=3,7.
\end{equation}
We denote the three-form on $X$ descended from $\beta_0$ by
$\bar\beta_0$.  It satisfies
$d\bar\beta_0=\upsilon_3-\upsilon_7$.
At $t=1$ one has $h=1$ and $f=0$.  Substitution of the component
parametrizations gives
\[
 \int_{N_1^k}\beta_{\mathcal C}=-1,\qquad
 \int_{N_2^k}\beta_{\mathcal C}=+1,\qquad
 \int_{N_j^k}\varphi=2\quad(j=1,2).
\]
Set $\sigma_1=-1$ and $\sigma_2=1$.  Every element of $K_0$ preserves
$\varphi$, preserves the Cayley orientations of these components, and
permutes the two lifts of each $N_j$.  Hence, for every
$\gamma\in K_0$ and $j,k\in\{1,2\}$,
\[
 \bigl[(\gamma^*\beta_{\mathcal C})|_{N_j^k}\bigr]
 =\frac{\sigma_j}{2}[\varphi|_{N_j^k}]\in H^3(N_j^k).
\]
There are eight terms in the definition of $\beta_0$ and its coefficient
is $1/4$.  It follows that
\[
 [\beta_0|_{N_j^k}]=\sigma_j[\varphi|_{N_j^k}].
\]
Both sides are pullbacks of classes on the quotient component $N_j$.
Pullback along the finite covering of the component is injective over
$\R$, so
\begin{equation}\label{eq:beta0-restrictions}
 [\bar\beta_0|_{N_1}]=-[\varphi|_{N_1}],\qquad
 [\bar\beta_0|_{N_2}]= [\varphi|_{N_2}].
\end{equation}
After shrinking the already chosen neighborhoods of $N_3$ and $N_7$, we
choose invariant tubular neighborhoods $V_1,V_2$ so that
$V_1,V_2,V_3,V_7$ are pairwise disjoint.  On $V_j$, for $j=1,2$, the
closed form
$\bar\beta_0-\sigma_j\pi_j^*(\varphi|_{N_j})$ restricts trivially to the
zero section.  Since $V_j$ deformation retracts onto $N_j$, there is a
two-form $q_j$ on $V_j$ such that
\[
 dq_j=\bar\beta_0-\sigma_j\pi_j^*(\varphi|_{N_j}).
\]
We average $q_j$ over the finite stabilizer of $V_j$ and choose an
invariant cutoff $\chi_j$ supported in $V_j$ and equal to one on a smaller
tubular neighborhood.  We define
\[
 \alpha_{\mathcal C}
 =\bar\beta_0-d(\chi_1q_1+\chi_2q_2).
\]
It follows directly that
\begin{equation}\label{eq:cobordism-primitive}
 d\alpha_{\mathcal C}=\upsilon_3-\upsilon_7,\qquad
 \alpha_{\mathcal C}|_{V_1}=-\pi_1^*(\varphi|_{N_1}),\qquad
 \alpha_{\mathcal C}|_{V_2}= \pi_2^*(\varphi|_{N_2}).
\end{equation}
The form $\bar\beta_0$ is zero in a neighborhood of the $t=1/2$ slice,
whereas the two cutoff corrections are supported near $N_1$ and $N_2$ in
the $t=1$ slice.  Therefore $\alpha_{\mathcal C}$ is zero near the
$t=1/2$ slice.  This slice contains every $N_j$ with
$3\leq j\leq10$, so
$\alpha_{\mathcal C}$ vanishes near all of those components, not only near
$N_3\cup N_7$.
The two signs are the oriented intersection numbers
$\mathcal C\cdot N_1=-1$ and $\mathcal C\cdot N_2=+1$; they follow by
comparing
$(-\partial_t,\partial_{x_1},\partial_{x_2},\partial_{x_3})$ on
$T\mathcal C$ with the Cayley orientations of $N_1,N_2$.

We choose the radial exceptional Thom representatives used in \cite{MM}:
\begin{equation}\label{eq:radial-Thom}
 \theta_j^{\rm MM}
 =\frac1{\pi\ii}\,d\!\left(b(r^2)\,\partial\log r^2\right),
\end{equation}
where $r$ is the radial coordinate on $\C^2$ in the resolution model
$Y\to\C^2/\{\pm1\}$, and $\partial$ is the Dolbeault $(1,0)$ operator in
those complex coordinates.  The function $b:[0,\infty)\to\R$ is smooth
and nondecreasing, equals $-1$ for
$s\leq\epsilon^2$, and equals $0$ for $s\geq4\epsilon^2$.  The form
$\theta_j^{\rm MM}$ integrates to one on the normal fibre
$\C/\{\pm1\}$ of the exceptional divisor.  Near that divisor it is the
pullback of a two-form on $\CP^1$, so
$(\theta_j^{\rm MM})^2$ vanishes on a neighborhood of the divisor.  On the
complement of the divisor, we transport the square through the blowdown
diffeomorphism.  Since it vanishes near the divisor, the transported form
extends by zero across $N_j$ and defines the smooth compactly supported
form $r_*((\theta_j^{\rm MM})^2)$ on $X$.  We choose the radial cutoff in
\eqref{eq:radial-Thom} sufficiently small that this support is contained in
the already fixed $V_j$; the form $\upsilon_j$ is supported there as well.
Fibre integration and the self-intersection number of
$\mathcal O_{\CP^1}(-2)$ give
\begin{equation}\label{eq:radial-Thom-square}
 [r_*((\theta_j^{\rm MM})^2)]
 =-2\operatorname{Th}[N_j]\qquad\text{in }H_c^4(V_j).
\end{equation}
Equation \eqref{eq:half-Thom} gives the exact compactly supported
cohomology calculation in $H_c^4(V_j)$
\[
 [r_*((\theta_j^{\rm MM})^2)+4\upsilon_j]
 =-2\operatorname{Th}[N_j]+4\cdot\frac12\operatorname{Th}[N_j]=0.
\]
By compactly supported de~Rham theory on the orbifold chart, there are
three-forms $\mu_j\in\Omega_c^3(V_j)$, extended by zero to $X$, such that
\[
 d\mu_j=r_*((\theta_j^{\rm MM})^2)+4\upsilon_j,
 \qquad j=3,7.
\]
The transported square and the annular representative $\upsilon_j$ both
vanish on a sufficiently small neighborhood of $N_j$.  Thus $d\mu_j=0$
there.  We denote the radial tubular homotopy operator by $H_j$.  Its
homotopy identity reduces on this neighborhood to
\[
 \mu_j-\pi_j^*(\mu_j|_{N_j})=dH_j\mu_j.
\]
We choose a cutoff $\chi_j^{\rm cut}$ that is supported in this
neighborhood and equals one on a smaller neighborhood.  Replacing $\mu_j$
by
\[
 \mu_j-d(\chi_j^{\rm cut}H_j\mu_j)
\]
does not change $d\mu_j$ and gives
$\mu_j=\pi_j^*(\mu_j|_{N_j})$ on the smaller neighborhood.  We retain the
notation $\mu_j$ for the modified form.  The cutoff correction is supported
in $V_j$, so the modified form still belongs to $\Omega_c^3(V_j)$.  Its
normal-pullback form implies
that $r^*\mu_j$ extends smoothly across the exceptional divisor.  Likewise
$r^*\alpha_{\mathcal C}$ is smooth:
it is zero near every $N_j$, $3\le j\le10$, and is the pullback of a form
on the centre near $N_1,N_2$.  Finally,
\begin{equation}\label{eq:pull-push-square}
 r^*r_*((\theta_j^{\rm MM})^2)=(\theta_j^{\rm MM})^2,
\end{equation}
because both sides vanish near $E_j$ and $r$ is a diffeomorphism on their
remaining support.  Hence
\begin{equation}\label{eq:etaMM-construction}
 \eta_{\rm MM}=r^*(\mu_7+4\alpha_{\mathcal C}-\mu_3)
\end{equation}
is smooth.  The supports of $\mu_3,\mu_7$ are disjoint from the first two
resolution neighborhoods because the four tubular neighborhoods were
chosen pairwise disjoint.  Thus \eqref{eq:cobordism-primitive} gives
\[
 \eta_{\rm MM}|_{\operatorname{supp}\theta_1^{\rm MM}}
 =-4\operatorname{pr}_1^*(\varphi|_{N_1}),\qquad
 \eta_{\rm MM}|_{\operatorname{supp}\theta_2^{\rm MM}}
 = 4\operatorname{pr}_2^*(\varphi|_{N_2}).
\]
The defining equations for $\mu_3$, $\mu_7$, and
$\alpha_{\mathcal C}$ give
\begin{align*}
d\eta_{\rm MM}
 &=r^*\bigl(d\mu_7+4d\alpha_{\mathcal C}-d\mu_3\bigr)\\
 &=r^*\bigl(r_*((\theta_7^{\rm MM})^2)+4\upsilon_7
       +4(\upsilon_3-\upsilon_7)
       -r_*((\theta_3^{\rm MM})^2)-4\upsilon_3\bigr)\\
 &=(\theta_7^{\rm MM})^2-(\theta_3^{\rm MM})^2.
\end{align*}
The last equality follows from \eqref{eq:pull-push-square}.
The restrictions above and the supported Thom product formula then give
\begin{equation}\label{eq:MM-primitive-data}
[\eta_{\rm MM}\theta_1^{\rm MM}]=-4v_1,\qquad
[\eta_{\rm MM}\theta_2^{\rm MM}]=4v_2.
\end{equation}
Thus the explicitly computed class is
\begin{equation}\label{eq:YMM}
Y_{\rm MM}
=[\eta_{\rm MM}(\theta_1^{\rm MM}+\theta_2^{\rm MM})]
=-4v_1+4v_2.
\end{equation}
With the convention \eqref{eq:Massey-coset}, the corresponding
defining-system representative is $-Y_{\rm MM}$; this sign does not affect
membership in the vector subspace $I_M$.
The projection of the \emph{entire} degree-five indeterminacy
$aH^3(M)+H^3(M)c$ to the exceptional summands is contained in
\begin{equation}\label{eq:indet-projection}
\operatorname{span}\{v_1+v_2,v_7-v_3\}.
\end{equation}
Indeed, if $\xi\in H^3(X)$, then
$\xi x_j=(\xi|_{N_j})x_j$.  We now perform the restriction calculation
explicitly.  Every such class has a representative
\begin{equation}\label{eq:H3X-normal-form}
 \xi=dt\wedge\beta_2
 +\lambda_1(dz_{123}+d\bar z_{123})
 +\lambda_2\bigl(dz_1\wedge dz_2\wedge d\bar z_3
 +d\bar z_1\wedge d\bar z_2\wedge dz_3\bigr),
\end{equation}
where $\beta_2$ is a constant closed two-form on the flat cover satisfying
$F^*\beta_2=\beta_2$ and $\kappa^*\beta_2=-\beta_2$.
The $dt$ term restricts to zero.  Substitution of the component
parametrizations from \eqref{eq:component-lifts} gives, on their chosen
twofold covers,
\begin{align*}
 \xi|_{N_1^1}&=4(\lambda_1-\lambda_2)\,
 dx_1\wedge dx_2\wedge dy_3,&
 \varphi|_{N_1^1}&=2\,dx_1\wedge dx_2\wedge dy_3,\\
 \xi|_{N_2^1}&=-4(\lambda_1-\lambda_2)\,
 dx_1\wedge dx_2\wedge dy_3,&
 \varphi|_{N_2^1}&=-2\,dx_1\wedge dx_2\wedge dy_3,\\
 \xi|_{N_3^1}&=2(\lambda_1+\lambda_2)\,
 dx_1\wedge dx_2\wedge dx_3,&
 \varphi|_{N_3^1}&=dx_1\wedge dx_2\wedge dx_3,\\
 \xi|_{N_7^1}&=-2(\lambda_1+\lambda_2)\,
 dy_1\wedge dy_2\wedge dx_3,&
 \varphi|_{N_7^1}&=-dy_1\wedge dy_2\wedge dx_3.
\end{align*}
Passing through the twofold covers changes both integrals by the same
factor.  Since
$H^3(N_j)=\R[\varphi|_{N_j}]$, the scalars are therefore
\begin{equation}\label{eq:restriction-scalars}
 \lambda_{12}(\xi)=2(\lambda_1-\lambda_2),\qquad
 \lambda_{37}(\xi)=2(\lambda_1+\lambda_2),
\end{equation}
and
\begin{align*}
[\xi|_{N_1}]&=\lambda_{12}(\xi)[\varphi|_{N_1}],&
[\xi|_{N_2}]&=\lambda_{12}(\xi)[\varphi|_{N_2}],\\
[\xi|_{N_3}]&=\lambda_{37}(\xi)[\varphi|_{N_3}],&
[\xi|_{N_7}]&=\lambda_{37}(\xi)[\varphi|_{N_7}].
\end{align*}
Consequently
\begin{equation}\label{eq:base-indet-projection}
\pi_E(a\xi)=\lambda_{12}(\xi)(v_1+v_2),\qquad
\pi_E(\xi c)=\lambda_{37}(\xi)(v_7-v_3).
\end{equation}
For the remaining terms, decompose an arbitrary class as
$h=\xi+\sum_j\lambda_jx_j\in H^3(M)$, with
$\xi\in H^3(X)$ and $\lambda_j\in H^1(N_j)$.  The supported representative
product rule of \cite[Section 3]{MM} is
\begin{equation}\label{eq:pushforward-product}
(\lambda x_j)x_j=-2(i_j)_!\lambda\in H^5(X),
\qquad i_j:N_j\hookrightarrow X,
\end{equation}
while $(\lambda x_j)x_k=0$ for $j\ne k$.  Hence these products have zero
exceptional projection.  More explicitly, multiplication by
$a=x_1+x_2$ can involve only the exceptional terms with $j=1,2$, and
multiplication by $c=x_7-x_3$ can involve only those with $j=3,7$.
Formula \eqref{eq:pushforward-product} places every resulting square in
the base summand $H^5(X)$, so none contributes to $\pi_E$.  Together with
\eqref{eq:base-indet-projection}, this proves
\eqref{eq:indet-projection}.  Finally, if
\[
-4v_1+4v_2=A(v_1+v_2)+B(v_7-v_3),
\]
the independent $v_1$ and $v_2$ coordinates in \eqref{eq:MM-cohom} would
give simultaneously $A=-4$ and $A=4$, which is impossible.  Therefore
$Y_{\rm MM}\notin I_M$, and likewise $-Y_{\rm MM}\notin I_M$.  Hence
\[
 0\notin\langle a,b,c\rangle.
\]
\end{proof}

Lemma~\ref{lem:tau-lift} gives
\begin{equation}\label{eq:parities}
\widetilde\tau^*a=-a,\qquad \widetilde\tau^*b=b,\qquad
\widetilde\tau^*c=c.
\end{equation}
We choose pairwise disjoint tubular neighborhoods of the exceptional
divisors that are invariant as a collection under $\widetilde\tau$.  We
choose a normalized Thom form $\theta_2$ representing $x_2$ and define
$\theta_1=-\widetilde\tau^*\theta_2$; equation \eqref{eq:tau-x} shows that
$\theta_1$ represents $x_1$.  For each $j\in\{3,7\}$, we average a
normalized Thom form with its pullback under $\widetilde\tau$.  The
resulting form, still denoted by $\theta_j$, is normalized and represents
$x_j$.  These forms satisfy
\begin{equation}\label{eq:thom-parities}
\widetilde\tau^*\theta_1=-\theta_2,\qquad
\widetilde\tau^*\theta_2=-\theta_1,\qquad
\widetilde\tau^*\theta_j=\theta_j\ (j=3,7).
\end{equation}
We define
\begin{equation}\label{eq:forms-abc}
\alpha=\theta_1+\theta_2,\qquad \beta=\theta_7+\theta_3,\qquad
\gamma=\theta_7-\theta_3.
\end{equation}
The disjoint supports give $\alpha\beta=0$, while
$\beta\gamma=\theta_7^2-\theta_3^2$ is exact by
$\PD[N_7]=\PD[N_3]$.  We choose a three-form $\eta_0$ satisfying
$d\eta_0=\beta\gamma$ and replace it by
\begin{equation}\label{eq:eta-average}
\eta=\tfrac12(\eta_0+\widetilde\tau^*\eta_0).
\end{equation}
Since $\beta\gamma$ is invariant, $d\eta=\beta\gamma$ and
$\widetilde\tau^*\eta=\eta$.  Moreover
$d(\alpha\eta)=\alpha\beta\gamma=0$.  Under the convention
\eqref{eq:Massey-coset}, the defining system
$(\xi_{12},\xi_{23})=(0,\eta)$ has representative $-\alpha\eta$.
Replacing the normalized Thom representatives or replacing a primitive by
another one changes this class by
$aH^3(M)+H^3(M)c=I_M$, exactly as in
\eqref{eq:general-indeterminacy}.  Thus the averaged system and the
explicit system \eqref{eq:MM-primitive-data} define representatives in the
same affine Massey coset.  Equivalently, their cohomology classes differ by
an element of $I_M$.
Proposition~\ref{prop:MM-Massey} therefore gives
$-[\alpha\eta]\notin I_M$.  Since $I_M$ is a vector subspace, we define
\begin{equation}\label{eq:Y-def}Y=[\alpha\eta]\in H^5(M).\end{equation}
Then $Y\notin I_M$, and
$\widetilde\tau^*Y=[(-\alpha)\eta]=-Y$.

\subsection{The quotient CDGA and its indeterminacy}
\begin{lemma}\label{lem:finite-quotient-CDGA}
Suppose that a finite group $\Gamma$ acts smoothly on a compact manifold
$P$, and denote the quotient map by $q:P\to P/\Gamma$.  With real coefficients,
pullback is a quasi-isomorphism
\begin{equation}\label{eq:APL-invariants}
q^*:A_{\rm PL}(P/\Gamma;\R)\longrightarrow
A_{\rm PL}(P;\R)^\Gamma.
\end{equation}
\end{lemma}
\begin{proof}
By the equivariant triangulation theorem for smooth finite-group actions
\cite{Illman},
$P$ admits a finite $\Gamma$-CW structure.  The image in $P/\Gamma$ of an
equivariant cell $\Gamma/H\times D^n$ is the ordinary cell $D^n$, whereas
the corresponding piece of the Borel construction is $BH\times D^n$.
The map from the Borel construction to the orbit space restricts over this
cell to the projection $BH\times D^n\to D^n$.  Since $H$ is finite,
$H^q(BH;\R)=0$ for $q>0$ and $H^0(BH;\R)=\R$.  Induction over the relative
cohomology sequences of the cellular skeleta shows that
$P_{h\Gamma}\to P/\Gamma$ is a real cohomology equivalence.  Equivalently,
the Leray spectral sequence has higher stalk at the orbit of $p$ equal to
$H^q(B\Gamma_p;\R)=0$ for every $q>0$.

Let $\pi_h:P_{h\Gamma}\to P/\Gamma$ be the preceding map.  Fix
$e\in E\Gamma$ and define $j_e:P\to P_{h\Gamma}$ by
$j_e(p)=[e,p]$.  Then $\pi_h\circ j_e=q$.  The Borel spectral sequence has
\[
E_2^{r,s}=H^r(\Gamma;H^s(P;\R)).
\]
Higher finite-group cohomology vanishes in characteristic zero, so it
collapses, and its edge homomorphism $j_e^*$ is an isomorphism from
$H^*(P_{h\Gamma};\R)$ onto $H^*(P;\R)^\Gamma$.  Since $\pi_h^*$ is an
isomorphism by the preceding skeletal argument and
$q^*=j_e^*\pi_h^*$, pullback gives
\[
q^*:H^*(P/\Gamma;\R)\xrightarrow{\ \cong\ }H^*(P;\R)^\Gamma.
\]
The Reynolds operator makes invariants exact, whence
\[
H^*(A_{\rm PL}(P;\R)^\Gamma)=H^*(A_{\rm PL}(P;\R))^\Gamma.
\]
This proves \eqref{eq:APL-invariants}.  The simplicial de~Rham theorem
\cite{BG,Dupont}, applied to the
smooth singular simplicial set and followed by exact finite-group
invariants, identifies this model with $\Omega(P)^\Gamma$ through a
zigzag of CDGA quasi-isomorphisms.

Finally, if $f:A\to B$ is a quasi-isomorphism of CDGAs, then for every
defined triple product
\[
H(f)\bigl(\langle a_1,a_2,a_3\rangle_A\bigr)
=\langle H(f)a_1,H(f)a_2,H(f)a_3\rangle_B.
\]
Indeed, $f$ carries one defining system to a defining system.  Both sides
are nonempty affine cosets, and $H(f)$ maps the entire indeterminacy
isomorphically by \eqref{eq:general-indeterminacy}; hence the initial
inclusion is equality.  Thus the two full Massey sets correspond under
$H(f)$.
\end{proof}

We denote the quotient map by $q:W=S^1\times M\to Z_0$.  The preceding lemma and
the equivariant simplicial de~Rham comparison give the explicit zigzag
\begin{equation}\label{eq:quotient-zigzag}
 A_{\rm PL}(Z_0;\R)\xrightarrow{\ q^*\ }
 A_{\rm PL}(W;\R)^{\langle\sigma\rangle}
 \xrightarrow{\ \simeq\ }
 A_{\rm dR}(\operatorname{Sing}^{\infty}W)^{\langle\sigma\rangle}
 \xleftarrow{\ \simeq\ }\Omega(W)^{\langle\sigma\rangle}.
\end{equation}
Here $A_{\rm dR}(\operatorname{Sing}^{\infty}W)$ is the CDGA of compatible
smooth forms on smooth singular simplices.  The left comparison includes
polynomial simplex forms and the right comparison pulls a differential
form back to every smooth simplex.  Both are natural
quasi-isomorphisms; taking invariants preserves this because the Reynolds
operator is exact.
Thus differential forms below live in the right-hand invariant de~Rham
model; their classes and their entire Massey cosets are transported to
$A_{\rm PL}(Z_0;\R)$ through \eqref{eq:quotient-zigzag}.  No de~Rham complex
on the singular quotient is being assumed.  More precisely, when an
invariant differential-form defining system is used below, equality of
Massey sets under the quasi-isomorphisms in the zigzag places the class of
its representative in the corresponding Massey set in
$A_{\rm PL}(Z_0;\R)$.  By the definition of that set, we then choose an
actual $A_{\rm PL}(Z_0;\R)$ defining system realizing the transported
class.

We write $u=[ds]\in H^1(S^1)$.  This class does not itself descend.  If
$\xi\in H^*(M)^-$, we write $u\xi$ for the unique quotient class whose
pullback to $S^1\times M$ is $[ds]\otimes\xi$.  The invariant ring
decomposition is
\begin{equation}\label{eq:quotient-cohom}
H^k(Z_0)=H^k(M)^+\oplus uH^{k-1}(M)^-.
\end{equation}
Here the superscripts denote the $\pm1$ eigenspaces of
$\widetilde\tau^*$.
Multiplication is inherited from $H^*(S^1\times M)$, in particular $u^2=0$.
The invariant forms on $W$
\begin{equation}\label{eq:ABC-forms}
\mathsf A=ds\wedge\alpha,\qquad \mathsf B=\beta,\qquad \mathsf C=\gamma
\end{equation}
have degrees $(3,2,2)$, with $\mathsf A\mathsf B=0$ and
$\mathsf B\mathsf C=d\eta$.  Thus the triple product
$\langle ua,b,c\rangle$ is defined.  The defining system
$(\xi_{12},\xi_{23})=(0,\eta)$ has representative
\[
0\cdot\mathsf C-(-1)^3\mathsf A\eta
=ds\wedge\alpha\wedge\eta,
\]
whose cohomology class corresponds to $uY\in H^6(Z_0)$ under
\eqref{eq:quotient-zigzag}.  Consequently there are closed elements
$\mathsf A_0,\mathsf B_0,\mathsf C_0\in A_{\rm PL}(Z_0;\R)$ representing
$ua,b,c$, respectively, and elements $\xi_{12}^0,\xi_{23}^0$ such that
\begin{align}
 d\xi_{12}^0&=\mathsf A_0\mathsf B_0,&
 d\xi_{23}^0&=\mathsf B_0\mathsf C_0,
 \label{eq:APL-defining-system}\\
 \left[\xi_{12}^0\mathsf C_0-(-1)^3
 \mathsf A_0\xi_{23}^0\right]&=uY\in H^6(Z_0;\R).
 \label{eq:APL-representative}
\end{align}
We fix this $A_{\rm PL}$ defining system for the remainder of the proof.

\begin{proposition}\label{prop:quotient-Massey}
The triple Massey product $\langle ua,b,c\rangle$ is defined, and
\begin{equation}\label{eq:quotient-nonzero}0\notin\langle ua,b,c\rangle.\end{equation}
\end{proposition}
\begin{proof}
In the two degrees entering the indeterminacy,
\[
H^3(Z_0)=H^3(M)^+\oplus uH^2(M)^-,\qquad
H^4(Z_0)=H^4(M)^+\oplus uH^3(M)^-.
\]
Since $u^2=0$, the two summands multiply as
\[
(ua)H^3(Z_0)=u\,aH^3(M)^+,\qquad
H^4(Z_0)c=H^4(M)^+c\oplus uH^3(M)^-c.
\]
The full indeterminacy is
\begin{align}
I_0&=(ua)H^3(Z_0)+H^4(Z_0)c\label{eq:I0-first}\\
&=H^4(M)^+c+u\bigl(aH^3(M)^++H^3(M)^-c\bigr).\label{eq:I0-decomp}
\end{align}
If $uY\in I_0$, projection to the $u$-summand would give
$Y\in aH^3(M)^++H^3(M)^-c\subset I_M$, a contradiction.  Thus the affine
coset $uY+I_0$ does not contain zero.
\end{proof}

\subsection{Passage through the resolution}
\begin{lemma}\label{lem:Gysin}
The space $Z_0$ is a compact oriented real homology eight-manifold, and the
degree-one map $\rho:\widehat Z\to Z_0$ has homomorphisms
\[
\rho_!:H^k(\widehat Z;\R)\longrightarrow H^k(Z_0;\R)
\qquad(0\leq k\leq8)
\]
such that, for every $v\in H^p(Z_0;\R)$ and
$w\in H^q(\widehat Z;\R)$,
\begin{equation}\label{eq:Gysin-properties}
\rho_!\rho^*v=v,\qquad
\rho_!(\rho^*v\smile w)=v\smile\rho_!(w).
\end{equation}
\end{lemma}
\begin{proof}
On the dense regular set $\widetilde\tau$ agrees with $\tau$, whose linear
part has determinant $-1$, so $\widetilde\tau$ reverses orientation.
Reflection reverses the orientation of $S^1$; hence $\sigma$ preserves the
orientation of $W$.

We fix $z\in Z_0$ and denote the stabilizer of one of its lifts by
$H\le\langle\sigma\rangle$.  The exponential map of an $H$-invariant metric gives the usual
Bochner linearization, so a small link of $z$ is $S^7/H$, with $H$ acting
linearly through $SO(8)$.  If $H$ is nontrivial, the derivative of its
generator is $+I$ on the tangent four-plane to the fixed component and
$-I$ on its normal four-plane; in particular its determinant and its degree
on $S^7$ are $+1$.  We choose an $H$-equivariant triangulation
\cite{Illman} and barycentrically subdivide it so that every simplex
stabilizer fixes that simplex pointwise.  The cellular chain complex of
the quotient is the coinvariant complex.  Over $\R$, the coinvariant
functor is exact, and the averaging operator identifies coinvariants with
invariants.  Therefore
\[
 H_*(S^7/H;\R)\cong H_*(S^7;\R)_H
 \cong H_*(S^7;\R)^H.
\]
Every element of $H$ has degree $+1$ on $S^7$, so this vector space is
$\R$ in degrees $0,7$ and zero otherwise.  Therefore
\begin{equation}\label{eq:local-homology}
H_i(Z_0,Z_0\setminus\{z\};\R)
\cong\widetilde H_{i-1}(S^7/H;\R)
=\begin{cases}\R,&i=8,\\0,&i\ne8.\end{cases}
\end{equation}
The local top generators agree with the quotient orientation.  Thus $Z_0$
is a compact oriented real homology eight-manifold.  Finite smooth
quotients are triangulable and locally contractible, so Poincar\'e duality
for compact oriented ENR homology manifolds
\cite[Chapter V]{BredonSheaf} applies: $Z_0$ has a
fundamental class $[Z_0]$ and cap product induces isomorphisms
\begin{equation}\label{eq:PD-homology-manifold}
 D_{Z_0}:H^k(Z_0)\longrightarrow H_{8-k}(Z_0),\qquad
 D_{Z_0}(v)=v\frown[Z_0].
\end{equation}

We orient $\widehat Z$ by its Cayley orientation.  We choose a point in
the nonempty common flat open set lying outside every first- and second-stage
resolution neighborhood.  At that point $\rho$ is the identity under the
staged/Joyce atlas identification.  The quotient orientation of $Z_0$
there is the orientation induced by the flat Cayley form $\Phi_0$; the
Joyce preglued positive form and its sufficiently small torsion-free
perturbation lie in the same positive Cayley cone and induce that same
orientation.  This regular point therefore has exactly one preimage and
local degree $+1$.
Hence $\rho_*[\widehat Z]=[Z_0]$.  We define
\begin{equation}\label{eq:Gysin-definition}
 \rho_!=D_{Z_0}^{-1}\rho_*D_{\widehat Z}.
\end{equation}
Cap-product naturality gives the projection formula
$\rho_!(\rho^*v\smile w)=v\smile\rho_!(w)$.  The fundamental-class
identity gives $\rho_!(1)=1$; therefore
\[
 \rho_!\rho^*v=\rho_!(\rho^*v\smile1)
 =v\smile\rho_!(1)=v,
\]
which proves both identities in \eqref{eq:Gysin-properties}.
\end{proof}

\begin{theorem}\label{thm:nonformal}
The triple Massey product
$\langle\rho^*(ua),\rho^*b,\rho^*c\rangle$ is defined, and
\begin{equation}\label{eq:final-Massey}
0\notin\langle\rho^*(ua),\rho^*b,\rho^*c\rangle.
\end{equation}
Thus $\widehat Z$ is nonformal.
\end{theorem}
\begin{proof}
The continuous map $\rho$ induces the strict CDGA morphism
\[
 \rho^*:A_{\rm PL}(Z_0;\R)\longrightarrow
 A_{\rm PL}(\widehat Z;\R).
\]
Pulling back the defining system fixed in
\eqref{eq:APL-defining-system}--\eqref{eq:APL-representative} proves that
the triple product is defined and gives a representative with cohomology
class $\rho^*(uY)$.  The ordinary de~Rham comparison for the smooth
manifold $\widehat Z$ carries this entire Massey set to the corresponding
one in $\Omega(\widehat Z)$.  Its full indeterminacy is
\begin{equation}\label{eq:Ihat}
I_{\widehat Z}=\rho^*(ua)H^3(\widehat Z)+H^4(\widehat Z)\rho^*c.
\end{equation}
Suppose
$\rho^*(uY)=\rho^*(ua)r+q\rho^*c$ for
$r\in H^3(\widehat Z)$ and $q\in H^4(\widehat Z)$.  The projection formula
and $\rho_!\rho^*=\id$ give
\[
uY=(ua)\rho_!(r)+\rho_!(q)c.
\]
For the second term, explicitly,
$\rho_!(q\rho^*c)=\rho_!(\rho^*c\,q)=c\,\rho_!(q)=\rho_!(q)c$; no sign
occurs because $|q||c|=8$.  The displayed class lies in $I_0$, contrary to
Proposition~\ref{prop:quotient-Massey}.  Hence
$\rho^*(uY)\notin I_{\widehat Z}$, so the full affine Massey coset does not
contain zero.  Every defined Massey product on a formal CDGA
contains zero, so \eqref{eq:final-Massey} obstructs formality of
$\widehat Z$.
\end{proof}

\section{Simple connectivity and exact holonomy}
\begin{proposition}\label{prop:pi1}
The manifold $\widehat Z$ is simply connected.
\end{proposition}
\begin{proof}
The manifold $M$ is simply connected \cite[Section 2.2]{MM}; hence
$\R\times M$ is simply connected.  We define an action on $\R\times M$
of the infinite dihedral group
\begin{equation}\label{eq:Dinf}
D_\infty=\langle S,R\mid R^2=1,\ RSR=S^{-1}\rangle,
\end{equation}
by the formulas $S(s,p)=(s+1,p)$ and
$R(s,p)=(-s,\widetilde\tau p)$.  These formulas satisfy the displayed
relations and therefore define the asserted action.  The fixed set of
$\widetilde\tau$ is nonempty.  Indeed, on $\T^7$ we impose
$t=\frac12$ and $z_3=\frac18$.  The resulting
$\tau$-fixed slice is a four-torus.  The lift of the old singular set is a
finite union of three-tori, so it cannot cover the slice.  We choose
$(z_1,z_2)\in E^2$ in its complement.  The resulting point is fixed by $\tau$.  Its
image $x\in X$ is regular; because $r$ is a diffeomorphism over the regular
set, the unique point of $r^{-1}(x)$ is fixed by $\widetilde\tau$.
Consequently $R$ fixes $(0,p)$ and $SR$ fixes $(\frac12,p)$ in
$\R\times M$.  The quotient by $\langle S\rangle$ is $S^1\times M$, and
the residual action of $R$ is precisely $\sigma$; hence
$(\R\times M)/D_\infty\cong Z_0$.  The space $\R\times M$ is connected,
simply connected, locally compact, and metrizable.  For every compact
subset $C\subset\R\times M$, its projection to the $s$-coordinate is
contained in a bounded interval.  Consequently only finitely many
integers $n$ satisfy $S^nC\cap C\ne\varnothing$, and only finitely many
integers $n$ satisfy $S^nRC\cap C\ne\varnothing$.
No nonidentity translation $S^n$ has a fixed point.  A reflection
$S^nR$ can fix $(s,p)$ only if $s=n/2$ and
$\widetilde\tau(p)=p$.  If two reflections $S^nR$ and $S^mR$ fixed the
same point, then their product $S^{n-m}$ would fix that point, and hence
$n=m$.  Thus every stabilizer is either trivial or the order-two group
generated by a unique reflection.

We also verify Armstrong's local discontinuity condition.  Given a point,
choose a relatively compact neighborhood.  The two compact-intersection
statements above leave only finitely many group elements whose translates
can meet its closure.  For each of those finitely many elements that does
not stabilize the point, shrink the neighborhood so that it is disjoint
from its translate.  The resulting neighborhood can meet a translate of
itself only under an element of its finite stabilizer.  Thus the action is
discontinuous in the precise sense required by
Armstrong's theorem \cite{Armstrong}, which gives
\[
\pi_1(Z_0)=D_\infty/N,
\]
where $N$ is normally generated by elements having fixed points.
Thus $R,SR\in N$, and $S=(SR)R\in N$, proving $\pi_1(Z_0)=1$.

We denote by $Z_{\rm ext}$ the complement of the interiors of mutually
disjoint conical neighborhoods of the singular components, and we set
$F=\Fix(\sigma)\subset W$.  Since $F$ has codimension four in the connected
manifold $W$, ordinary relative transversality shows that $W\setminus F$
is path connected.  Its quotient is the regular locus of $Z_0$.  Radial
deformation in each punctured cone fibre moves paths out of the deleted
tubular interiors and onto their boundaries; hence $Z_{\rm ext}$ is path
connected.  We choose a basepoint in $Z_{\rm ext}$ and compatible
basepaths to the boundary components.  For one component $B$ of $F$, which
we identify with its image in $Z_0$, we denote
the boundary $\RP^3$-bundle by $L\to B$, the cone bundle by $P_0\to B$,
and the truncated Eguchi--Hanson bundle by $P_1\to B$; the latter is the
bundle in \eqref{eq:associated-EH-bundle}.  Their fibres are
$\RP^3$, $C(\RP^3)$, and the disc bundle of
$\mathcal O_{\CP^1}(-2)$, respectively.  The last two fibres are connected
and simply connected.  The homotopy exact sequences therefore give
isomorphisms
\[
(p_i)_*:\pi_1(P_i)\xrightarrow{\ \cong\ }\pi_1(B),\qquad i=0,1.
\]
If $j_i:L\to P_i$ are the boundary inclusions and $p_L:L\to B$ is the
projection, then $(p_i)_*(j_i)_*=(p_L)_*$.  After identifying both
$\pi_1(P_i)$ with $\pi_1(B)$ by $(p_i)_*$, the two boundary homomorphisms
are therefore the same homomorphism $(p_L)_*$.  This argument neither assumes that
$\pi_1(\RP^3)\to\pi_1(L)$ is injective nor ignores monodromy.  Van Kampen
applies to open collar thickenings of $Z_{\rm ext}$ and $P_i$, whose
intersection deformation retracts onto $L$, and gives the pushouts
\begin{equation}\label{eq:pi1-pushouts}
 \pi_1(Z_{\rm ext}\cup_LP_i)
 \cong \pi_1(Z_{\rm ext})*_{\pi_1(L)}\pi_1(P_i),
 \qquad i=0,1,
\end{equation}
with identical structure maps after the preceding identifications.  The
fibrewise construction in Lemma~\ref{lem:two-stage} glues $P_1$ to
$Z_{\rm ext}$ by the same fixed boundary identification used to glue
$P_0$.  Therefore the two groups in
\eqref{eq:pi1-pushouts} are isomorphic through the specified basepoint,
basepath, and bundle-projection identifications.

A compact smooth fixed set has finitely many components.  Enumerate them
as $B_1,\ldots,B_k$.  For $0\leq j\leq k$, let $Z^{(j)}$ be the space in
which the cone bundles over $B_1,\ldots,B_j$ have been replaced by their
truncated Eguchi--Hanson bundles, while the cone bundles over the remaining
components have not been replaced.  Choose all basepaths in the common
exterior.  Applying the one-component pushout calculation to
$Z^{(j-1)}$ and $Z^{(j)}$ gives an isomorphism
$\pi_1(Z^{(j-1)})\cong\pi_1(Z^{(j)})$ through those specified
identifications.  Induction on $j$ therefore gives
$\pi_1(\widehat Z)\cong\pi_1(Z_0)=1$.
\end{proof}

\begin{theorem}\label{thm:full-holonomy}
For every sufficiently small $t>0$, the metric
$g_t=g_{\widehat\Phi_t}$ has full Riemannian holonomy
$\Hol(\widehat Z,g_t)=\SpinSeven$.
\end{theorem}
\begin{proof}
We fix a sufficiently small $t$ and set $g_t=g_{\widehat\Phi_t}$.  A
torsion-free $\SpinSeven$ structure makes $g_t$ Ricci-flat and determines a
parallel unit positive spinor.
Since $\widehat Z$ is compact, $g_t$ is complete.  Since $\widehat Z$ is
simply connected, full and restricted holonomy agree.  Moreover,
$\pi_1(\widehat Z)=0$ implies $H_1(\widehat Z;\Z)=0$ by abelianization,
and the universal coefficient theorem gives
$H^1(\widehat Z;\R)=0$.  A parallel vector would dualize to a nonzero
parallel one-form.  Such a form is harmonic and represents a nonzero
cohomology class, which is impossible.  Hence the de Rham decomposition
has no Euclidean factor.

The global de Rham decomposition theorem applies because
$(\widehat Z,g_t)$ is complete and simply connected.  Since there is no
Euclidean factor, it gives a Riemannian product
\[
 (\widehat Z,g_t)=(X_1,g_1)\times\cdots\times(X_r,g_r)
\]
of irreducible, complete, simply connected, nonflat factors.  Each $X_i$
is spin: restricting $w_2(T\widehat Z)=0$ to $X_i$ times basepoints in the
other factors gives $w_2(TX_i)=0$.  We denote the restricted
holonomy group of $X_i$ by $H_i\subset SO(n_i)$.  The factor spin
structure determines its connected spin-holonomy lift
$\widetilde H_i\subset\operatorname{Spin}(n_i)$.

Under the standard homomorphism
$\operatorname{Spin}(n_1)\times\cdots\times\operatorname{Spin}(n_r)
\to\operatorname{Spin}(8)$, the restriction of
the complex positive spin module is a finite direct sum of external
tensor products of complex factor spin modules.  More explicitly, an
odd-dimensional factor contributes its irreducible complex spin module,
an even-dimensional factor contributes one of its two half-spin modules,
and the parity constraint selects the summands; when two odd-dimensional
factors are combined, the usual two parity summands occur.  This is the
standard branching rule for complex Clifford modules, and every projection
onto one of the displayed tensor summands is equivariant for
$\widetilde H_1\times\cdots\times\widetilde H_r$.
Therefore the parallel positive spinor has a nonzero invariant component
in at least one tensor summand.  If $\Delta_i$ denotes the factor spin or
half-spin module occurring in that summand, then
\[
 (\Delta_1\otimes\cdots\otimes\Delta_r)
 ^{\widetilde H_1\times\cdots\times\widetilde H_r}
 =\Delta_1^{\widetilde H_1}\otimes\cdots\otimes
   \Delta_r^{\widetilde H_r}.
\]
The nonzero invariant component therefore implies
$\Delta_i^{\widetilde H_i}\ne0$ for every $i$.  By the holonomy
principle, every factor $X_i$ therefore admits a nonzero parallel spinor.

The Berger--Simons--Wang classification \cite{Wang} now applies to each
factor.  An irreducible nonsymmetric factor with a parallel spinor has
holonomy $SU(m)$, $Sp(k)$, $G_2$, or $\SpinSeven$ in real dimension $2m$,
$4k$, $7$, or $8$, respectively.  An irreducible locally symmetric space
is Einstein with nonzero Einstein constant unless it is Euclidean.
Consequently an irreducible Ricci-flat symmetric factor is flat and has
already been excluded.  The remaining dimension possibilities are
\[
\begin{array}{c|c|c}
\text{factor dimension}&\text{nonflat parallel-spinor holonomy}
 &\text{required complementary dimension}\\
\hline
8&\SpinSeven,\ SU(4),\ Sp(2)&0\\
7&G_2&1\text{ (flat)}\\
6&SU(3)&2\text{ (flat)}\\
4&Sp(1)&4.
\end{array}
\]
The Wang list has no nonflat irreducible factor in dimension five, and
Ricci-flat factors of dimensions one, two, or three are flat.
Consequently a six- or seven-dimensional factor forces a flat
two- or one-dimensional complement.  Without a flat factor, the only possibilities are
one eight-dimensional factor with holonomy $\SpinSeven$, $SU(4)$, or
$Sp(2)$, or two four-dimensional factors with holonomy
$Sp(1)\times Sp(1)$.  The de Rham product above is global, and full
holonomy equals restricted holonomy because $\widehat Z$ is simply
connected.  Thus no factor-permuting or other disconnected extension of
these connected holonomy groups occurs.

The projection of the compact product $\widehat Z$ onto each de Rham
factor is surjective, so every factor is compact.
Holonomy $SU(4)$ is K\"ahler, and holonomy $Sp(2)$ is hyperk\"ahler.  In
the $Sp(1)\times Sp(1)$ case, de Rham gives a product of two compact
hyperk\"ahler four-manifolds, which is again K\"ahler.  Every compact
K\"ahler manifold is formal by \cite{DGMS}, contradicting
Theorem~\ref{thm:nonformal}.  Therefore
$\Hol(\widehat Z,g_t)=\SpinSeven$.
\end{proof}

Under the identification of Lemma~\ref{lem:two-stage}, the map appearing
in Theorem~\ref{thm:main} is
\[
 \pi=q_0\circ\rho:\widehat Z\longrightarrow\mathcal O.
\]
It is the same map as the Joyce blowdown
\eqref{eq:Joyce-blowdown}.  Its properness was proved immediately after
\eqref{eq:Joyce-blowdown}, and the local model is the identity over the
regular stratum, so it is a diffeomorphism over the regular locus.
Let $\varepsilon_{\rm J}$ be the upper bound supplied by
Theorem~\ref{thm:Joyce-existence}, and let $\varepsilon_{\rm H}$ be the
smallness bound in Theorem~\ref{thm:full-holonomy}.  Taking
$\varepsilon=\min\{\varepsilon_{\rm J},\varepsilon_{\rm H}\}$ makes the
existence and full-holonomy conclusions hold for every
$0<t<\varepsilon$.  Theorem~\ref{thm:nonformal} and
Proposition~\ref{prop:pi1} supply the remaining two assertions.  This
proves Theorem~\ref{thm:main}.

\appendix
\section{Lean formalization of the algebraic deductions}\label{app:lean}

The formalization encodes the algebraic implications used in the Massey,
fundamental-group, and holonomy arguments.  The geometric statements proved
in the body of the paper enter through explicitly named hypotheses.

\subsection{Affine cosets and the exceptional-coordinate obstruction}

Let $V$ be an abelian group, let $I\leq V$, and let
\[
 \mathcal C(y,I)=y+I=\{y+i:i\in I\}.
\]
The definitions \texttt{AdditiveCore}, \texttt{AdditiveSubset}, and
\texttt{affineCoset} encode these data.  The theorem
\texttt{zero\_mem\_affineCoset\_iff} proves
\begin{equation}\label{eq:lean-coset}
 0\in\mathcal C(y,I)\quad\Longleftrightarrow\quad y\in I.
\end{equation}
Indeed, an equality $0=y+i$ with $i\in I$ gives $y=-i\in I$, and the
converse follows from $0=y+(-y)$.  Consequently, the theorem
\texttt{full\_affine\_coset\_exclusion} proves
\[
 y\notin I\quad\Longrightarrow\quad 0\notin\mathcal C(y,I).
\]
If $y'\in y+I$ and $y\notin I$, then $y'\notin I$ as well.  This implication
is the theorem
\[
 \text{\texttt{translated\_representative\_exclusion}}.
\]

For the exceptional projection in Section~5, write the coordinates in the
order $(v_1,v_2,v_3,v_7)$.  The projected indeterminacy is
\begin{equation}\label{eq:lean-coordinate-subspace}
 J=\{(A,A,-B,B):A,B\in\R\},
\end{equation}
whereas the projected representative is
\[
 y_{\rm MM}=(-4,4,0,0).
\]
If $y_{\rm MM}\in J$, its first two coordinates give $A=-4$ and $A=4$.
Thus $y_{\rm MM}\notin J$.  The theorem
\texttt{unequal\_coordinate\_certificate} proves this implication for an
arbitrary coefficient type from the inequality of the two coordinates,
and \texttt{yMM\_not\_in\_projected\_indeterminacy} gives the corresponding
decidable calculation over $\Z$.  The theorem
\texttt{exclusion\_by\_projection} then lifts the coordinate obstruction
from $J$ to the full indeterminacy subgroup.

\subsection{Passage through the quotient and the resolution}

The quotient step is expressed by groups $V,Q$, maps
$\iota_u:V\to Q$ and $p_u:Q\to V$, and indeterminacy subgroups
$I_V\leq V$ and $I_Q\leq Q$.  The hypotheses are
\begin{equation}\label{eq:lean-quotient-hypotheses}
 p_u\iota_u=\id_V,\qquad p_u(I_Q)\subseteq I_V.
\end{equation}
The theorem \texttt{quotient\_indeterminacy\_transfer} proves
\begin{equation}\label{eq:lean-quotient-transfer}
 y\notin I_V\quad\Longrightarrow\quad \iota_u(y)\notin I_Q.
\end{equation}
This is the formal counterpart of projection to the $u$-summand in
Proposition~\ref{prop:quotient-Massey}.

For the resolution step, let $\rho^*:Q\to R$ and $\rho_!:R\to Q$ satisfy
\begin{equation}\label{eq:lean-gysin-hypotheses}
 \rho_!\rho^*=\id_Q,\qquad \rho_!(I_R)\subseteq I_Q.
\end{equation}
The structure \texttt{GysinFormulaData} records separately the two terms
in the triple-product indeterminacy and the corresponding projection
formulas.  The theorems \texttt{gysin\_maps\_full\_indeterminacy} and
\texttt{gysin\_indeterminacy\_transfer} prove
\begin{equation}\label{eq:lean-gysin-transfer}
 q\notin I_Q\quad\Longrightarrow\quad \rho^*q\notin I_R.
\end{equation}
The structure \texttt{MasseyPipelineInput} combines
\eqref{eq:lean-coset}--\eqref{eq:lean-gysin-transfer}.  Its conclusion is
the theorem
\[
 \text{\texttt{full\_massey\_coset\_excludes\_zero}},
\]
which states that zero is absent from the full Massey coset on the resolved
space.

\subsection{The dihedral and holonomy deductions}

The infinite dihedral group is represented in normal form by
\[
 D_\infty=\{S^n:n\in\Z\}\sqcup\{S^nR:n\in\Z\},
 \qquad R^2=1,\qquad RSR=S^{-1}.
\]
If a normal subgroup $N\triangleleft D_\infty$ contains $R$ and $SR$, then
$S=(SR)R\in N$.  Hence $S^n\in N$ and $S^nR\in N$ for every $n\in\Z$, so
$N=D_\infty$.  This is the theorem \texttt{DInf.normal\_generation}; the
structure \texttt{ArmstrongInput} combines it with the hypotheses of
Armstrong's theorem and the invariance of the fundamental group under the
fibrewise resolution.

The toric calculation is encoded using doubled support values.  From
$0<\ell$ and $2\ell<L$, the theorem
\texttt{toric\_support\_sign\_certificate} proves
\[
 -(0+0-2\ell)=2\ell>0,
 \qquad
 -(\ell+\ell-L)=L-2\ell>0.
\]
These are the two forms of wall pairing needed in Section~3.2.

Finally, the formalized holonomy alternative is
\[
 \SpinSeven,\qquad SU(4),\qquad Sp(2),\qquad Sp(1)\times Sp(1).
\]
The three proper alternatives imply that the manifold is K\"ahler, and
the DGMS theorem then implies formality.  Under the hypothesis of
nonformality, \texttt{holonomy\_alternative\_elimination} therefore selects
$\SpinSeven$.

\subsection{Assembly}

The structure \texttt{GeometricInput} contains the following propositions:
the resolution is smooth and compact; its blowdown is proper and is a
diffeomorphism over the regular locus; it carries a torsion-free
$\SpinSeven$ structure; the cohomological data satisfy the Massey pipeline;
the covering action satisfies the hypotheses of Armstrong's theorem; the
fibrewise resolution preserves the fundamental group; the holonomy belongs
to the four alternatives above; the three proper alternatives are
K\"ahler; compact K\"ahler manifolds are formal; and formality forces zero
to belong to every defined Massey product.  These propositions are exactly
the geometric, topological, and cohomological results established in the
body or supplied by the cited theorems.

For every value $D$ of \texttt{GeometricInput}, the theorem
\texttt{verified\_conclusion} produces the conjunction that the resolved
space is smooth, compact, and simply connected; the blowdown is proper and
is a diffeomorphism over the regular locus; the torsion-free structure has
full holonomy $\SpinSeven$; and the full Massey coset excludes zero, so the
resolved space is nonformal.  This is the formalized logical assembly of
Theorem~\ref{thm:main}.

\begin{thebibliography}{99}
\bibitem{Armstrong}M.~A. Armstrong,
\emph{The fundamental group of the orbit space of a discontinuous group},
Proc. Cambridge Philos. Soc. \textbf{64} (1968), 299--301.
\bibitem{BG}A.~K. Bousfield and V.~K.~A.~M. Gugenheim,
\emph{On PL de Rham theory and rational homotopy type},
Mem. Amer. Math. Soc. \textbf{8} (1976), no.~179.
\bibitem{BredonSheaf}G.~E. Bredon,
\emph{Sheaf Theory}, second ed., Graduate Texts in Mathematics 170,
Springer, New York, 1997.
\bibitem{BredonGroup}G.~E. Bredon,
\emph{Introduction to Compact Transformation Groups}, Pure and Applied
Mathematics 46, Academic Press, New York--London, 1972.
\bibitem{DGMS}P.~Deligne, P.~Griffiths, J.~Morgan, and D.~Sullivan,
\emph{Real homotopy theory of K\"ahler manifolds},
Invent. Math. \textbf{29} (1975), 245--274.
\bibitem{Dupont}J.~L. Dupont,
\emph{Curvature and Characteristic Classes}, Lecture Notes in Mathematics
640, Springer, Berlin, 1978.
\bibitem{Fulton}W.~Fulton,
\emph{Introduction to Toric Varieties}, Annals of Mathematics Studies 131,
Princeton University Press, Princeton, 1993.
\bibitem{Illman}S.~Illman,
\emph{The equivariant triangulation theorem for actions of compact Lie
groups}, Math. Ann. \textbf{262} (1983), 487--501.
\bibitem{Hirsch}M.~W. Hirsch,
\emph{Differential Topology}, Graduate Texts in Mathematics 33,
Springer, New York, 1976.
\bibitem{JoyceBook}D.~D. Joyce,
\emph{Compact Manifolds with Special Holonomy}, Oxford Mathematical
Monographs, Oxford University Press, 2000.
\bibitem{MM}L.~Mart\'in--Merch\'an,
\emph{Compact holonomy $G_2$ manifolds need not be formal},
Geom. Topol. \textbf{30} (2026), 1109--1127;
doi:10.2140/gt.2026.30.1109.
\bibitem{Wang}M.~Y. Wang,
\emph{Parallel spinors and parallel forms},
Ann. Global Anal. Geom. \textbf{7} (1989), 59--68;
doi:10.1007/BF00137402.
\end{thebibliography}
\end{document}
